\section{Reduction of symmetric sets to case classes}\label{sr:sec}

This section gives written proofs of the declarations of \dirn{SYMLEAN/SymmetricReduction.lean}. Its main result,
Theorem~\ref{sr:thm:main} (\lean{heptagon_of_nontrivial_isometry}), is a conditional statement. Five propositions
about point sets of an exact size with one prescribed symmetry, the \emph{case classes} \lean{HalfTurn34},
\lean{Threefold33}, \lean{Fivefold35}, \lean{MirrorOneAxis33} and \lean{MirrorNoAxis34}
(Definition~\ref{sr:def:classes}), imply the following: every finite set of at least $33$ points in general position
that is mapped into itself by a Euclidean isometry of the plane other than the identity contains $7$ points in convex
position. The classes are \emph{hypotheses} of the theorem. Nothing in this section proves them;
Remark~\ref{sr:rem:records} says which project records supply them and with which rigor label. The section also proves
the variant Theorem~\ref{sr:thm:variant} (\lean{halfTurn_of_symmetric_counterexample}), in which a fourfold class
\lean{Fourfold32} replaces \lean{HalfTurn34}, and the bridge Lemma~\ref{sr:lem:mirror-line} to a mirror statement of
the directory \dirn{H3}.

The proof of Theorem~\ref{sr:thm:main} has three parts.
\begin{enumerate}
\item Classification (Theorem~\ref{sr:thm:structure}). A Euclidean isometry $g\ne\mathrm{id}$ of the plane that maps
  a finite set with at least two points into itself is a rotation $z\mapsto c+a(z-c)$ with $a\ne1$ a root of unity, or a
  reflection in a line. Translations and glide reflections are excluded because the set is finite.
\item Rotations. A power of $g$ is the half-turn, or a rotation whose factor is a primitive $p$-th root of unity for an
  odd prime $p$ (Lemma~\ref{sr:lem:root-cases}). Whole orbits are selected (Lemma~\ref{sr:lem:orbit}) to reach the
  exact sizes of the classes: $17$ antipodal pairs; $11$ orbits of size $3$; $7$ orbits of size $5$. For $p\ge7$ one
  orbit consists of $p$ concyclic points, which are in convex position (Lemma~\ref{sr:lem:concyclic}).
\item Reflections. At most two points lie on the mirror and the others come in pairs; $17$ pairs, or $16$ pairs and one
  point of the mirror, are kept.
\end{enumerate}

\paragraph{Sources.}
Paths of research records are relative to \dirn{cloud/es7-2026-10-19/results}; paths of the main library are relative
to the repository root. The first 12 hex digits of the SHA-256 of each Lean source used are:
\begin{itemize}
\item \dirn{SYMLEAN/SymmetricReduction.lean} (1059 lines; every declaration of this file is covered here):
  \texttt{6d5ec90a4215};
\item \dirn{SYMLEAN/SymmetricReductionAxioms.lean}: \texttt{7bf809cbdc77}. Its lines 1--1059 are byte-identical to
  the previous file; two \texttt{\#print axioms} commands follow;
\item the eight repository modules that the file imports, directly or through other imports:
  \begin{itemize}
  \item \dirn{ErdosSzekeres/Definitions.lean}: \texttt{2425500eefd1};
  \item \dirn{ErdosSzekeres/Basic.lean}: \texttt{13667aa2d3dc};
  \item \dirn{ErdosSzekeres/Orientation.lean}: \texttt{f67955ac7091};
  \item \dirn{ErdosSzekeres/OrientationGeometry.lean}: \texttt{c7db9bbbb65e};
  \item \dirn{ErdosSzekeres/Suk/SupportingLines.lean}: \texttt{6a8e12677685};
  \item \dirn{ErdosSzekeres/Suk/BaseOrder.lean}: \texttt{d21becd29e84};
  \item \dirn{ErdosSzekeres/Suk/BaseChart.lean}: \texttt{08d0e856f005};
  \item \dirn{ErdosSzekeres/Suk/TriangleOrder.lean}: \texttt{6447cd4833d4}.
  \end{itemize}
  The last three are imported only through \dirn{ErdosSzekeres/Suk/SupportingLines.lean}, and this section cites
  nothing from them.
\end{itemize}
The files were checked with Lean 4.24.0 and Mathlib \texttt{f897ebcf}.

\paragraph{Axioms and rigor.}
\begin{itemize}
\item The compile log \dirn{SYMLEAN/logs/SymmetricReduction.out} is empty, and the exit code is 0
  (\dirn{SYMLEAN/logs/runs.tsv}). So no declaration of the file uses \lean{sorry}.
\item The log \dirn{SYMLEAN/logs/SymmetricReductionAxioms.out} shows that the two theorems
  \lean{heptagon_of_nontrivial_isometry} and \lean{halfTurn_of_symmetric_counterexample} depend only on the axioms
  \lean{propext}, \lean{Classical.choice} and \lean{Quot.sound}.
\item The bridge lemma \lean{mirrorOneAxis33_of_exclusionLine} is used in the proof of the review theorem
  \lean{Review.mirrorOneAxis33_iff_exclusionLine} at \dirn{SYMLEAN/review/ReviewSanity.lean:1629}. That file's lines 1--1059
  are byte-identical to \dirn{SYMLEAN/SymmetricReduction.lean}, and its log \dirn{SYMLEAN/logs/ReviewSanity.out} shows
  the same three axioms for this theorem.
\item The file and the eight imported repository modules listed above contain no \lean{axiom} declaration, no
  \lean{sorry}, no \lean{decide} and no \lean{native_decide} (the word ``axiom'' occurs there only in comments). A
  declaration of the file that the printed theorems do not use (such as
  \lean{IsEucIsometry.injective}) is covered by the empty compile log and by this scan: it uses no \lean{sorry}, and
  neither the file nor these modules add an axiom. No statement of this section is a finite computation.
\end{itemize}

Every statement below whose head carries a Lean reference is \emph{Lean-checked} in this sense. The proofs written here
are \emph{written proofs (not refereed)}. Unless a proof says otherwise, it follows the Lean proof step by step.
Statements tagged ``written step'' are introduced only for the written proofs; most of them are standard facts that
the Lean proofs take from Mathlib, and each is proved here.

\paragraph{Dependencies.}
The section uses the foundations: Definitions~\ref{fd:def:point}, \ref{fd:def:orient}, \ref{fd:def:gp},
\ref{fd:def:cp}, \ref{fd:def:ngon}, \ref{fd:def:inv} and \ref{fd:def:rot}, and Lemmas~\ref{fd:lem:collinear},
\ref{fd:lem:cp-hull} and~\ref{fd:lem:mono}. Lemma~\ref{sr:lem:supports} is the lemma titled ``Strict supporting
lines'' in the foundations section (the same declaration of the main library), repeated here with its short proof.
Remarks~\ref{sr:rem:records}
and~\ref{sr:rem:main} refer to Theorems~\ref{cb:thm:c2} and~\ref{cb:thm:c3} (bookkeeping of two case splits; both
theorems are computer-assisted, as their Lean proofs rest on kernel evaluation, and Theorem~\ref{cb:thm:c2} is
Lean-checked according to the record, Theorem~\ref{cb:thm:c3} Lean-checked), to Theorem~\ref{mb:thm:classes}, which
gives written proofs of the two mirror classes, to Theorem~\ref{mb:thm:main3} and to Theorem~\ref{ar:thm:main}; no
proof of this section uses them.

\subsection{Setting and definitions}\label{sr:ssec:defs}

\paragraph{Conventions.}
Points are elements of $\R^2$, written $p=(p_0,p_1)$ (Definition~\ref{fd:def:point}; Lean type \lean{Point}). Point
sets are finite (Lean \texttt{Finset Point}), and $\card S$ is the number of elements of $S$. The function $\orient$ is
that of Definition~\ref{fd:def:orient}: $\orient(a,b,c)=(b_0-a_0)(c_1-a_1)-(b_1-a_1)(c_0-a_0)$. General position,
convex position and \texttt{HasConvexNGon n S} are as in Definitions~\ref{fd:def:gp}, \ref{fd:def:cp}
and~\ref{fd:def:ngon}. We use them in the following forms.
\begin{itemize}
\item $S$ is in general position if and only if $\orient(a,b,c)\ne0$ for all pairwise distinct $a,b,c\in S$
  (Lemma~\ref{fd:lem:collinear}).
\item $S$ is in convex position if and only if $p\notin\conv(S\setminus\{p\})$ for every $p\in S$
  (Lemma~\ref{fd:lem:cp-hull}).
\item \texttt{HasConvexNGon n S} means: some $T\subseteq S$ with $\card T=n$ is in convex position.
\item Monotonicity (Lemma~\ref{fd:lem:mono}): subsets of a set in general position are in general position; subsets of
  a set in convex position are in convex position; \texttt{HasConvexNGon n T} and $T\subseteq S$ imply
  \texttt{HasConvexNGon n S}.
\item If $T$ is in convex position and $n\le\card T$, then \texttt{HasConvexNGon n T}: any $n$-element subset of $T$ is
  in convex position by monotonicity (Lean \lean{ConvexPosition.hasConvexNGon}, \dirn{ErdosSzekeres/Basic.lean:31}).
\end{itemize}
For a map $f\colon X\to X$ the iterates are $f^{[0]}=\mathrm{id}$ and $f^{[k+1]}=f\circ f^{[k]}$; then also
$f^{[k+1]}=f^{[k]}\circ f$ and $f^{[i+j]}=f^{[i]}\circ f^{[j]}$. Natural numbers include $0$, and every natural number
divides $0$. For $z\in\C$, $\bar z$ is the complex conjugate (Mathlib \lean{starRingEnd}) and $\abs z$ the modulus;
Mathlib's $\mathrm{normSq}(z)=(\operatorname{Re}z)^2+(\operatorname{Im}z)^2$ equals $\abs z^2$ and is multiplicative.
We write $\langle z,w\rangle=\operatorname{Re}(z\bar w)=\operatorname{Re}z\operatorname{Re}w+\operatorname{Im}z
\operatorname{Im}w$ for the Euclidean inner product of $\C=\R^2$, so that $\abs{z}^2=\langle z,z\rangle$.

Lean names in the text are read relative to the namespace \lean{ErdosSzekeres.SymReduction}: a name $n$ stands for the
first of \lean{ErdosSzekeres.SymReduction.}$n$, \lean{ErdosSzekeres.}$n$ and $n$ that exists. So \lean{orbit_select}
is \lean{ErdosSzekeres.SymReduction.orbit_select}, \lean{Review.} and \lean{VariantChecks.} denote sub-namespaces of
\lean{ErdosSzekeres.SymReduction}, \lean{Point} is \lean{ErdosSzekeres.Point}, \lean{Suk.} denotes
\lean{ErdosSzekeres.Suk}, and the names of Mathlib and those with the prefixes \lean{Mir.} (the Lean package in
\dirn{H3}) and \lean{V2L.} (the Lean package in \dirn{V2-L}) are full names. Statement heads give full names.

\begin{definition}[Squared distance, Euclidean isometry; \leanref{ErdosSzekeres.SymReduction.eDist2}{SYMLEAN/SymmetricReduction.lean:265}, \leanref{ErdosSzekeres.SymReduction.IsEucIsometry}{SYMLEAN/SymmetricReduction.lean:268}]\label{sr:def:iso}
For $p,q\in\R^2$ put $\operatorname{eDist2}(p,q)=(p_0-q_0)^2+(p_1-q_1)^2$. A map $g\colon\R^2\to\R^2$ is a
\emph{Euclidean isometry} (Lean \texttt{IsEucIsometry g}) if
\[
  \operatorname{eDist2}(g(p),g(q))=\operatorname{eDist2}(p,q)\qquad\text{for all }p,q\in\R^2 .
\]
\end{definition}

No further property of $g$ is assumed: $g$ is not assumed continuous, surjective or affine, but it is defined on the
whole plane. That $g$, read in complex coordinates, has the form $z\mapsto az+b$ or $z\mapsto a\bar z+b$ with
$\abs a=1$ follows from Lemma~\ref{sr:lem:toC-iso}(b) and Theorem~\ref{sr:thm:structure}(a).

\begin{remark}[Why an explicit Euclidean distance]\label{sr:rem:sup}
Mathlib equips the type \lean{Point} of functions $\{0,1\}\to\R$ with the product (sup) metric
$\mathrm{dist}(p,q)=\max(\abs{p_0-q_0},\abs{p_1-q_1})$; for instance $\mathrm{dist}((0,0),(3,4))=4$. Mathlib's predicate
\lean{Isometry} for maps of \lean{Point} therefore means preservation of the sup distance, which a rotation need not
have. For example, the rotation $\rho$ by $2\pi/3$ about a point $c$ fixes $c$ and maps $c+(1,0)$, at sup distance $1$
from $c$, to $c+(-\tfrac12,\tfrac{\sqrt3}{2})$, at sup distance $\tfrac{\sqrt3}{2}$ from $c$. So $\rho$ does not preserve
the sup distance, and a theorem stated with Mathlib's \lean{Isometry} would not cover the rotations used in
\lean{Threefold33}. This is why the file states isometries with the explicit squared Euclidean distance
$\operatorname{eDist2}$. General position, convex position and \lean{HasConvexNGon} are affine notions (Mathlib
\lean{Collinear} and \lean{ConvexIndependent}) and involve no metric. Lean checks of such sup-metric facts are in the
review file \dirn{SYMLEAN/review/SupMetric.lean}, which is outside the scope of this section.
\end{remark}

\begin{definition}[Invariance; \leanref{ErdosSzekeres.SymReduction.Invariant}{SYMLEAN/SymmetricReduction.lean:271}]\label{sr:def:inv}
For $f\colon\R^2\to\R^2$ and a finite set $S$, $\operatorname{Invariant}(f,S)$ means $f(p)\in S$ for every $p\in S$.
\end{definition}

This is the notion of Definition~\ref{fd:def:inv}: $f$ maps $S$ \emph{into} $S$; equality $f(S)=S$ is not required.
(It follows when $f$ is injective, because $S$ is finite, but it is never used.)

\begin{definition}[Complex coordinates; \leanref{ErdosSzekeres.SymReduction.toC}{SYMLEAN/SymmetricReduction.lean:274}, \leanref{ErdosSzekeres.SymReduction.ofC}{SYMLEAN/SymmetricReduction.lean:277}]\label{sr:def:complex}
For $p\in\R^2$ let $\operatorname{toC}(p)=p_0+ip_1\in\C$, and for $z\in\C$ let $\operatorname{ofC}(z)=(\operatorname{Re}
z,\operatorname{Im}z)\in\R^2$. In proofs we abbreviate $\hat p=\operatorname{toC}(p)$.
\end{definition}

\begin{lemma}[Coordinates; \leanref{ErdosSzekeres.SymReduction.toC_re}{SYMLEAN/SymmetricReduction.lean:279}, \leanref{ErdosSzekeres.SymReduction.toC_im}{SYMLEAN/SymmetricReduction.lean:280}, \leanref{ErdosSzekeres.SymReduction.ofC_zero}{SYMLEAN/SymmetricReduction.lean:281}, \leanref{ErdosSzekeres.SymReduction.ofC_one}{SYMLEAN/SymmetricReduction.lean:282}, \leanref{ErdosSzekeres.SymReduction.toC_ofC}{SYMLEAN/SymmetricReduction.lean:284}, \leanref{ErdosSzekeres.SymReduction.ofC_toC}{SYMLEAN/SymmetricReduction.lean:287}, \leanref{ErdosSzekeres.SymReduction.toC_injective}{SYMLEAN/SymmetricReduction.lean:291}, \leanref{ErdosSzekeres.SymReduction.eDist2_eq_normSq}{SYMLEAN/SymmetricReduction.lean:294}, \leanref{ErdosSzekeres.SymReduction.eDist2_pos_of_ne}{SYMLEAN/SymmetricReduction.lean:297}, \leanref{ErdosSzekeres.SymReduction.point_ext}{SYMLEAN/SymmetricReduction.lean:502}]\label{sr:lem:complex}
Let $p,q\in\R^2$ and $z\in\C$.
\begin{enumerate}[label=(\alph*)]
\item (\lean{toC_re}, \lean{toC_im}, \lean{ofC_zero}, \lean{ofC_one}) $\operatorname{Re}\operatorname{toC}(p)=p_0$,
  $\operatorname{Im}\operatorname{toC}(p)=p_1$,
  $\operatorname{ofC}(z)_0=\operatorname{Re}z$ and $\operatorname{ofC}(z)_1=\operatorname{Im}z$.
\item (\lean{point_ext}) If $p_0=q_0$ and $p_1=q_1$, then $p=q$.
\item (\lean{toC_ofC}, \lean{ofC_toC}, \lean{toC_injective}) $\operatorname{toC}(\operatorname{ofC}(z))=z$ and
  $\operatorname{ofC}(\operatorname{toC}(p))=p$; the map
  $\operatorname{toC}$ is injective.
\item (\lean{eDist2_eq_normSq}) $\operatorname{eDist2}(p,q)=\mathrm{normSq}(\operatorname{toC}(p)-\operatorname{toC}(q))$.
\item (\lean{eDist2_pos_of_ne}) If $p\ne q$, then $0<\operatorname{eDist2}(p,q)$.
\end{enumerate}
\end{lemma}

\begin{proof}
(a) holds by the definitions. (b) A point is a function on $\{0,1\}$, determined by its two values. (c) By (a),
$\operatorname{toC}(\operatorname{ofC}(z))=\operatorname{Re}z+i\operatorname{Im}z=z$, and
$\operatorname{ofC}(\operatorname{toC}(p))$ has coordinates $p_0,p_1$, so it equals $p$ by (b). If
$\operatorname{toC}(p)=\operatorname{toC}(q)$, then $p=\operatorname{ofC}(\operatorname{toC}(p))=
\operatorname{ofC}(\operatorname{toC}(q))=q$. (d) $\operatorname{toC}(p)-\operatorname{toC}(q)=(p_0-q_0)+i(p_1-q_1)$,
whose $\mathrm{normSq}$ is $(p_0-q_0)^2+(p_1-q_1)^2$. (e) If $p\ne q$, then $\hat p-\hat q\ne0$ by (c), and a nonzero
complex number has positive $\mathrm{normSq}$; conclude by (d).
\end{proof}

\begin{definition}[Rotations, half-turn, line reflection; \leanref{ErdosSzekeres.SymReduction.crot}{SYMLEAN/SymmetricReduction.lean:349}, \leanref{ErdosSzekeres.SymReduction.rotAbout}{SYMLEAN/SymmetricReduction.lean:400}, \leanref{ErdosSzekeres.SymReduction.halfTurn}{SYMLEAN/SymmetricReduction.lean:413}, \leanref{ErdosSzekeres.SymReduction.reflLine}{SYMLEAN/SymmetricReduction.lean:423}]\label{sr:def:maps}
Let $p\in\R^2$.
\begin{enumerate}[label=(\alph*)]
\item For $c,w\in\C$ (Lean \texttt{crot c w}):
  $\operatorname{crot}_{c,w}(p)=\operatorname{ofC}\bigl(c+w(\operatorname{toC}(p)-c)\bigr)$.
\item For $c\in\R^2$ and $\theta\in\R$ (Lean \lean{rotAbout}; the formula of Definition~\ref{fd:def:rot}):
  \[
  \operatorname{rotAbout}_{c,\theta}(p)=\bigl(c_0+\cos\theta\,(p_0-c_0)-\sin\theta\,(p_1-c_1),\;
  c_1+\sin\theta\,(p_0-c_0)+\cos\theta\,(p_1-c_1)\bigr).
  \]
\item For $c\in\R^2$ (Lean \texttt{halfTurn c}): $\operatorname{halfTurn}_c(p)=(2c_0-p_0,\,2c_1-p_1)$.
\item For $a,u\in\R^2$ (Lean \texttt{reflLine a u}): with $d=(p_0-a_0)u_0+(p_1-a_1)u_1$,
  \[
  \operatorname{reflLine}_{a,u}(p)=\bigl(a_0+2du_0-(p_0-a_0),\;a_1+2du_1-(p_1-a_1)\bigr).
  \]
\end{enumerate}
\end{definition}

For $\abs w=1$, $\operatorname{crot}_{c,w}$ is the rotation about the point $\operatorname{ofC}(c)$ through the angle
$\theta$ with $w=e^{i\theta}$ (Lemma~\ref{sr:lem:rot-half}(a)), and $\operatorname{rotAbout}_{c,\theta}$ is the
counterclockwise rotation by $\theta$ about $c$. For $u_0^2+u_1^2=1$, $\operatorname{reflLine}_{a,u}$ is the reflection
in the line $a+\R u$: by Lemma~\ref{sr:lem:refl}(a) it is $z\mapsto\hat a+\hat u^{\,2}\,\overline{(z-\hat a)}$ in
complex coordinates, and as $\hat u\,\overline{\hat u}=1$, this map fixes $\hat a+t\hat u$ and sends $\hat a+it\hat u$
to $\hat a-it\hat u$ for every $t\in\R$. For $u=0$ it is the half-turn about $a$. For every other $u$ it is not even a
Euclidean isometry: with $s=u_0^2+u_1^2\notin\{0,1\}$, it fixes $a$ and sends $a+u$ to $a+(2s-1)u$, and
$\operatorname{eDist2}(a+(2s-1)u,a)=(2s-1)^2s\ne s=\operatorname{eDist2}(a+u,a)$. So every statement below that uses
it as a reflection assumes $u_0^2+u_1^2=1$. The definition has the same text as \lean{Mir.reflLine} in
\dirn{H3/MirrorEncoding/Line.lean:18}.

\begin{definition}[The case classes; \leanref{ErdosSzekeres.SymReduction.HalfTurn34}{SYMLEAN/SymmetricReduction.lean:467}, \leanref{ErdosSzekeres.SymReduction.Threefold33}{SYMLEAN/SymmetricReduction.lean:472}, \leanref{ErdosSzekeres.SymReduction.Fivefold35}{SYMLEAN/SymmetricReduction.lean:477}, \leanref{ErdosSzekeres.SymReduction.MirrorOneAxis33}{SYMLEAN/SymmetricReduction.lean:482}, \leanref{ErdosSzekeres.SymReduction.MirrorNoAxis34}{SYMLEAN/SymmetricReduction.lean:488}, \leanref{ErdosSzekeres.SymReduction.Fourfold32}{SYMLEAN/SymmetricReduction.lean:496}]\label{sr:def:classes}
The following six propositions are defined (Lean: definitions of type \lean{Prop}). In each, $S$ ranges over all finite
subsets of $\R^2$, and $c$, $a$, $u$ range over all points of $\R^2$.
\begin{enumerate}[label=(\arabic*)]
\item \lean{HalfTurn34}: for all $S$ and $c$, if $S$ is in general position, $\card S=34$, $c\notin S$ and
  $\operatorname{Invariant}(\operatorname{halfTurn}_c,S)$, then \texttt{HasConvexNGon 7 S}.
\item \lean{Threefold33}: for all $S$ and $c$, if $S$ is in general position, $\card S=33$, $c\notin S$ and
  $\operatorname{Invariant}(\operatorname{rotAbout}_{c,2\pi/3},S)$, then \texttt{HasConvexNGon 7 S}.
\item \lean{Fivefold35}: for all $S$ and $c$, if $S$ is in general position, $\card S=35$, $c\notin S$ and
  $\operatorname{Invariant}(\operatorname{rotAbout}_{c,2\pi/5},S)$, then \texttt{HasConvexNGon 7 S}.
\item \lean{MirrorOneAxis33}: for all $S$, $a$ and $u$, if $u_0^2+u_1^2=1$, $S$ is in general position, $\card S=33$,
  $\operatorname{Invariant}(\operatorname{reflLine}_{a,u},S)$, and there is $x\in S$ with
  $\operatorname{reflLine}_{a,u}(x)=x$ such that every $y\in S$ with $\operatorname{reflLine}_{a,u}(y)=y$ equals $x$, then
  \texttt{HasConvexNGon 7 S}.
\item \lean{MirrorNoAxis34}: for all $S$, $a$ and $u$, if $u_0^2+u_1^2=1$, $S$ is in general position, $\card S=34$,
  $\operatorname{Invariant}(\operatorname{reflLine}_{a,u},S)$, and $\operatorname{reflLine}_{a,u}(p)\ne p$ for every
  $p\in S$, then \texttt{HasConvexNGon 7 S}.
\item \lean{Fourfold32}: for all $S$ and $c$, if $S$ is in general position, $\card S=32$, $c\notin S$ and
  $\operatorname{Invariant}(\operatorname{rotAbout}_{c,2\pi/4},S)$, then \texttt{HasConvexNGon 7 S}.
\end{enumerate}
The angles are the real numbers $2\pi/3$, $2\pi/5$ and $2\pi/4$ (Lean \texttt{2 * Real.pi / 3} etc.); the rotations are
counterclockwise. In \lean{MirrorOneAxis33} the last hypothesis says that exactly one point of $S$ is fixed by the
reflection; in \lean{MirrorNoAxis34} no point of $S$ is fixed. For $u_0^2+u_1^2=1$ the fixed points of
$\operatorname{reflLine}_{a,u}$ are exactly the points of the mirror $a+\R u$. Indeed, with $d$ as in
Definition~\ref{sr:def:maps}(d), the equation $\operatorname{reflLine}_{a,u}(p)=p$ says $2du_0=2(p_0-a_0)$ and
$2du_1=2(p_1-a_1)$, so a fixed point is $p=a+du$; and for $p=a+tu$ one has $d=t(u_0^2+u_1^2)=t$ and
$\operatorname{reflLine}_{a,u}(p)=a+2tu-tu=p$.
\end{definition}

These are definitions, not assertions. No class is proved in this section. Theorem~\ref{sr:thm:main} takes the first
five as hypotheses, and Theorem~\ref{sr:thm:variant} takes \lean{Threefold33}, \lean{Fourfold32}, \lean{Fivefold35} and
the two mirror classes. Remark~\ref{sr:rem:redundant} records which hypotheses of the classes are implied by the others.

\begin{remark}[Where the project supplies the classes]\label{sr:rem:records}
None of the following is checked in this section; it follows \dirn{C2ENC/class_table.md}, \S5 and \S9 of
\dirn{SYMLEAN/README.md}, \S3 of \dirn{WRITEUP/SYMMETRIC.md} and the records named below. For no class is the step from
a SAT certificate to the Lean class Lean-checked. On the certificate routes of \lean{HalfTurn34}, \lean{Threefold33},
\lean{Fivefold35} and \lean{Fourfold32}, and for the \dirn{SYMZ} certificates of the two mirror classes, that step rests
on written proofs (the soundness of the encoders and the lemmas named below) and on the review of the class
definitions in \dirn{SYMLEAN/README.md} \S6, which compared each class with its certificate by reading. On the route
through \dirn{H3} and \dirn{V2-L} to \lean{MirrorOneAxis33}, the soundness of the encodings and the reduction of
\lean{Mir.MirrorExclusionLine} to CNF facts are Lean-checked. The cube covers in that reduction are checked by kernel
evaluation (computer-assisted), with one exception: the cover of the top two levels of the hull-pair case split of
sub-case (ii), where the point of $S$ on the mirror lies in the convex hull of the other points
(\dirn{WRITEUP/MIRROR.md} \S4). That cover is itself given by $1+121$ of those CNF facts. What remains outside Lean on
this route is described in Remark~\ref{sr:rem:mirror-line}. The
project trusts one proof checker, \texttt{cake\_lpr}: where a certificate below is said to be accepted by both
checkers, \texttt{cake\_lpr} accepted it, and \texttt{lrat-check} is only a secondary check (\dirn{WRITEUP/SYMMETRIC.md}
\S3, revised on 6 November 2026). The trust statement in \S9 of \dirn{SYMLEAN/README.md}, which this remark follows on
other points, predates that revision and still lists both checkers as trusted.
\begin{itemize}
\item \lean{HalfTurn34}: claimed in \dirn{C2FINAL} as computer-assisted. It rests on SAT certificates (\dirn{C2C34}:
  hull cover and $120$ hexagonal hulls, proofs retained; \dirn{C2P}--\dirn{C2P5}: parallelogram subcases, most of which
  are verdicts of both checkers recorded at solve time for proofs that were then deleted; re-checks in
  \dirn{C2RECHECK}); on the written soundness proof of the encoder, \dirn{C2ENC/proof.md}; on the reductions A and B,
  which pair equivalent parallelogram cases and subcases (written proofs, reviewed in \dirn{C2P/review} and checked at
  the CNF level in \dirn{AUDIT3}); and on the bookkeeping of the subcase orbits, Lean-checked according to the record
  (\dirn{C2FINAL/README.md} \S6), by kernel evaluation, and so computer-assisted (Theorem~\ref{cb:thm:c2}; that the
  listed subcases are refuted rests on the records). A table of the evidence is in \dirn{C2ENC/class_table.md} \S1.
\item \lean{Threefold33}: claimed in the records as computer-assisted. It rests on \dirn{C3FINAL}, a composition of
  $116{,}883$ pieces whose proofs were not retained, so that their verdicts are recorded ones (a random sample of
  $2{,}120$ pieces was
  re-solved in \dirn{C3SPOT}, and the $204$ indispensable $\sigma$-cube pieces were re-solved with retained proofs in
  \dirn{C3PIECES}); on the cover of the $910$ roots, Lean-checked by kernel evaluation and so computer-assisted
  (Theorem~\ref{cb:thm:c3}; that each leaf names a record with recorded checker verdicts is checked outside Lean); on
  the written normal form \dirn{C3-S/soundness.md} and the written Lemma H6M (\dirn{C3H6M/soundness_h6m.md}, checked
  adversarially in \dirn{C3AUDIT2}); on the written exhaustiveness of the layer splits and $\sigma$-splits (Lemmas 1--4
  and the cover theorem of \dirn{C3-S/soundness.md}, its version for the reduced cut in
  \dirn{C3H6M/soundness_h6m.md}, and rules 2 and 3 of \dirn{C3U2/sigma_split.md}); on the encoding \dirn{C3a}; and on
  the written centre-perturbation lemma \dirn{C2ENC/perturbation.md}. The perturbation lemma is needed because the
  class admits centres that lie on a line through two points of different orbits, which the encoder excludes.
\item \lean{Fivefold35}: claimed in the records as computer-assisted (\dirn{SYM5/README.md}). It rests on \dirn{SYM5}:
  two normal forms of one encoder (an angle form and a layer form; not independent encodings), each with a retained
  proof that both checkers accepted; either proof would suffice for the class without the other, given the written
  proofs listed next. It rests further on written proofs of the encoder's soundness and of the two normal forms, and on
  the same perturbation lemma. An earlier computation reported by B.~Dumitru (arXiv:2512.24061, \S7), which has
  priority, also implies the fivefold bound; no proof files of that run are public (\dirn{SYM5/README.md}).
\item \lean{MirrorOneAxis33} and \lean{MirrorNoAxis34}: Lean-checked in \dirn{SYMLEAN/mirror}
  (\lean{mirrorOneAxis33_proved}, \lean{mirrorNoAxis34_proved}), with no SAT input, from the repository's formalization
  of Baek's theorem (J.~Baek, arXiv:2206.04260); written proofs are in Theorem~\ref{mb:thm:classes}. The file
  \dirn{SYMLEAN/mirror/ReductionThreeClasses.lean} does not import \dirn{SYMLEAN/SymmetricReduction.lean}. Its lines
  38--1084 are byte-identical to lines 12--1058 of \dirn{SYMLEAN/SymmetricReduction.lean}, so the classes defined in
  \dirn{SYMLEAN/mirror/ReductionThreeClasses.lean} are copies of those of Definition~\ref{sr:def:classes}, with the same
  full names and the same text. Both classes are also supplied by SAT certificates in \dirn{SYMZ} (Theorems 2a and 2b
  of \dirn{WRITEUP/SYMMETRIC.md}, claimed in the records as computer-assisted), and \lean{MirrorOneAxis33} also,
  conditionally, by the Lean chain of \dirn{H3} and \dirn{V2-L} (Remark~\ref{sr:rem:mirror-line}).
\item \lean{Fourfold32} (among the statements of this section, a hypothesis only of Theorem~\ref{sr:thm:variant} and
  of Lemma~\ref{sr:lem:fourfold}, which its proof uses): claimed in \dirn{C4CUT}, which states it without a
  rigor label. It rests on a final certificate accepted by both checkers, on six cuts checked by exact computation and
  on a written proof of the encoder's soundness. The statement is due to Subercaseaux, Mackey, Qian and Heule
  (arXiv:2506.00224, \S6.2).
\end{itemize}
\end{remark}

\subsection{Finite dynamics}\label{sr:ssec:dyn}

\begin{lemma}[Iterates stay in an invariant set; \leanref{ErdosSzekeres.SymReduction.iterate_mem}{SYMLEAN/SymmetricReduction.lean:69}, \leanref{ErdosSzekeres.SymReduction.Invariant.iterate}{SYMLEAN/SymmetricReduction.lean:309}]\label{sr:lem:iterate}
Iterates of a map stay in a set that the map sends into itself.
\begin{enumerate}[label=(\alph*)]
\item (\lean{iterate_mem}) Let $X$ be a set, $f\colon X\to X$, and $S\subseteq X$ finite with $f(y)\in S$ for every
  $y\in S$. If $x\in S$, then $f^{[i]}(x)\in S$ for every $i\in\N$.
\item (\lean{Invariant.iterate}) If $f\colon\R^2\to\R^2$, $S\subseteq\R^2$ is finite and $\operatorname{Invariant}(f,S)$,
  then $\operatorname{Invariant}(f^{[k]},S)$ for every $k\in\N$.
\end{enumerate}
\end{lemma}

\begin{proof}
(a) Induction on $i$: $f^{[0]}(x)=x\in S$, and $f^{[i+1]}(x)=f(f^{[i]}(x))\in S$ because $f^{[i]}(x)\in S$. (b) is (a)
applied to each $x\in S$.
\end{proof}

\begin{lemma}[Pigeonhole for iterates; written step]\label{sr:lem:pigeon}
Let $S\subseteq\R^2$ be finite, $f\colon\R^2\to\R^2$ with $\operatorname{Invariant}(f,S)$, and $q\in S$. Then there are
$i<j$ in $\N$ with $f^{[i]}(q)=f^{[j]}(q)$.
\end{lemma}

\begin{proof}
By Lemma~\ref{sr:lem:iterate}, the $\card S+1$ points $f^{[k]}(q)$, $0\le k\le\card S$, lie in $S$. By the pigeonhole
principle two of them with different indices coincide; call the smaller index $i$ and the larger $j$. (In Lean this is
the local statement \lean{pigeon} in the proof of \lean{isometry_structure}, obtained from Mathlib's
\lean{Set.Finite.exists_lt_map_eq_of_forall_mem}.)
\end{proof}

\begin{definition}[Minimal period; written step]\label{sr:def:minper}
Let $X$ be a set, $f\colon X\to X$ and $x\in X$. If $f^{[n]}(x)=x$ for some $n>0$, then $\operatorname{minimalPeriod}(f,x)$
is the least such $n$; otherwise $\operatorname{minimalPeriod}(f,x)=0$. (This is Mathlib's
\lean{Function.minimalPeriod}.)
\end{definition}

\begin{lemma}[Minimal periods; written step]\label{sr:lem:minper}
Let $f\colon X\to X$, $x\in X$ and $d=\operatorname{minimalPeriod}(f,x)$.
\begin{enumerate}[label=(MP\arabic*)]
\item If $f^{[n]}(x)=x$ for some $n\in\N$, then $d\mid n$. (Mathlib \lean{Function.IsPeriodicPt.minimalPeriod_dvd}.)
\item $d=1$ if and only if $f(x)=x$. (Mathlib \lean{Function.minimalPeriod_eq_one_iff_isFixedPt}.)
\item $f^{[d]}(x)=x$. (Mathlib \lean{Function.iterate_minimalPeriod}.)
\item If $i<j<d$, then $f^{[i]}(x)\ne f^{[j]}(x)$. (Mathlib \lean{Function.iterate_injOn_Iio_minimalPeriod}.)
\end{enumerate}
\end{lemma}

\begin{proof}
(MP3) If some $n>0$ has $f^{[n]}(x)=x$, this holds for the least one, $d$. Otherwise $d=0$ and $f^{[0]}(x)=x$.

(MP1) If $n=0$ there is nothing to prove. Let $n>0$. Then $d>0$ is the least positive $m$ with $f^{[m]}(x)=x$. By
induction on $t$, $f^{[td]}(x)=x$ for all $t\in\N$. Write $n=td+r$ with $0\le r<d$. Then
$x=f^{[n]}(x)=f^{[r]}(f^{[td]}(x))=f^{[r]}(x)$. If $r>0$, this contradicts the minimality of $d$; so $r=0$ and $d\mid n$.

(MP2) If $f(x)=x$, then $1$ is the least positive integer and $f^{[1]}(x)=x$, so $d=1$. Conversely, if $d=1\ne0$, then by
definition $f^{[1]}(x)=x$.

(MP4) Suppose $f^{[i]}(x)=f^{[j]}(x)$ with $i<j<d$. Then $d>0$, and applying $f^{[d-j]}$ to both sides gives
$f^{[d-j+i]}(x)=f^{[d]}(x)=x$ by (MP3). Since $j<d$ and $i<j$, we have $0<d-j+i<d$, which contradicts the minimality
of $d$.
\end{proof}

\begin{definition}[Forward orbit segment; \leanref{ErdosSzekeres.SymReduction.orbitFS}{SYMLEAN/SymmetricReduction.lean:65}]\label{sr:def:orbitFS}
For $f\colon X\to X$, $p\in\N$ and $x\in X$, let $\operatorname{orbitFS}(f,p,x)=\{f^{[i]}(x):0\le i<p\}$, a finite set
(Lean: the image of $\{0,\dots,p-1\}$ under $i\mapsto f^{[i]}(x)$; Lean needs decidable equality on $X$ to form it).
\end{definition}

\begin{lemma}[Orbit segments; \leanref{ErdosSzekeres.SymReduction.card_orbitFS_of_minimalPeriod}{SYMLEAN/SymmetricReduction.lean:111}, \leanref{ErdosSzekeres.SymReduction.mem_orbitFS_self}{SYMLEAN/SymmetricReduction.lean:90}, \leanref{ErdosSzekeres.SymReduction.orbitFS_subset}{SYMLEAN/SymmetricReduction.lean:93}, \leanref{ErdosSzekeres.SymReduction.orbitFS_invariant}{SYMLEAN/SymmetricReduction.lean:99}, \leanref{ErdosSzekeres.SymReduction.card_orbitFS}{SYMLEAN/SymmetricReduction.lean:76}]\label{sr:lem:orbitFS}
Let $X$ be a set, $f\colon X\to X$, $p\in\N$ and $x\in X$, and write $O=\operatorname{orbitFS}(f,p,x)$.
\begin{enumerate}[label=(\alph*)]
\item (\lean{mem_orbitFS_self}) If $0<p$, then $x\in O$.
\item (\lean{orbitFS_subset}) If $S\subseteq X$ is finite, $f(y)\in S$ for every $y\in S$, and $x\in S$, then
  $O\subseteq S$.
\item (\lean{orbitFS_invariant}) If $0<p$ and $f^{[p]}(x)=x$, then $f(y)\in O$ for every $y\in O$.
\item (\lean{card_orbitFS}) If $p$ is prime, $f^{[p]}(x)=x$ and $f(x)\ne x$, then $\card O=p$.
\item (\lean{card_orbitFS_of_minimalPeriod}) If $\operatorname{minimalPeriod}(f,x)=p$, then $\card O=p$.
\end{enumerate}
\end{lemma}

\begin{proof}
(a) $x=f^{[0]}(x)$ and $0<p$. (b) Each $f^{[i]}(x)$ lies in $S$ by Lemma~\ref{sr:lem:iterate}(a). (c) Let
$y=f^{[i]}(x)$ with $i<p$. Then $f(y)=f^{[i+1]}(x)$. If $i+1<p$, this lies in $O$. Otherwise $i+1=p$ and
$f(y)=f^{[p]}(x)=x\in O$ by (a). (e) By (MP4), $i\mapsto f^{[i]}(x)$ is injective on $\{0,\dots,p-1\}$, so $O$ has $p$
elements. (This also holds for $p=0$, when $O=\varnothing$.) (d) By (MP1), $d=\operatorname{minimalPeriod}(f,x)$ divides
the prime $p$, so $d=1$ or $d=p$. By (MP2), $d\ne1$ because $f(x)\ne x$. Hence $d=p$, and (e) applies.
\end{proof}

The next statement contains the induction shared by the two orbit selection lemmas. The Lean file carries out the same
induction twice, once in each lemma; we do it once.

\begin{lemma}[Orbit selection, common form; written step]\label{sr:lem:orbit-core}
Let $X$ be a set, $f\colon X\to X$ and $p\in\N$ with $0<p$. Let $S\subseteq X$ be finite, and assume: $f(y)\in S$ for
every $y\in S$; $f$ is injective on $S$; and every $y\in S$ satisfies $f^{[p]}(y)=y$ and
$\card{\operatorname{orbitFS}(f,p,y)}=p$. Then $p\mid\card S$, and for every $k\in\N$ with $kp\le\card S$ there is
$T\subseteq S$ with $f(y)\in T$ for every $y\in T$ and $\card T=kp$.
\end{lemma}

\begin{proof}
Strong induction on $\card S$. The hypotheses on single points pass to subsets of $S$, and so does injectivity.

If $S=\varnothing$, then $p\mid0$; and $kp\le0$ forces $kp=0$, so $T=\varnothing$ works.

Let $S\ne\varnothing$ and pick $x\in S$. Put $O=\operatorname{orbitFS}(f,p,x)$. By Lemma~\ref{sr:lem:orbitFS} and the
hypotheses, $O\subseteq S$, $\card O=p$, $x\in O$, and $f(y)\in O$ for every $y\in O$. Let $S'=S\setminus O$; then
$\card{S'}=\card S-p<\card S$.

\emph{$S'$ is mapped into itself.} Let $y\in S'$. Then $f(y)\in S$. Suppose $f(y)\in O$, say $f(y)=f^{[i]}(x)$ with
$i<p$. If $i=0$, then $f(y)=x=f^{[p]}(x)=f(f^{[p-1]}(x))$ (here $p\ge1$). Both $y$ and $f^{[p-1]}(x)$ lie in $S$
(Lemma~\ref{sr:lem:iterate}(a)), so injectivity on $S$ gives $y=f^{[p-1]}(x)\in O$. If $i\ge1$, then
$f(y)=f(f^{[i-1]}(x))$ and injectivity gives $y=f^{[i-1]}(x)\in O$. Both contradict $y\notin O$. Hence $f(y)\in S'$.

So $S'$ satisfies all the hypotheses, and by induction $p\mid\card{S'}$, and $S'$ has $f$-invariant subsets of every
size $k'p\le\card{S'}$. Since $\card S=\card{S'}+p$, we get $p\mid\card S$. Let $kp\le\card S$. If $k=0$, take
$T=\varnothing$. If $k=k'+1$, then $k'p\le\card S-p=\card{S'}$. Take an $f$-invariant $T'\subseteq S'$ with
$\card{T'}=k'p$, and put $T=T'\cup O\subseteq S$. The union of two sets that $f$ maps into themselves is mapped into
itself, and $T'\cap O=\varnothing$ because $T'\subseteq S\setminus O$. So $\card T=k'p+p=kp$.
\end{proof}

\begin{lemma}[Orbit selection; \leanref{ErdosSzekeres.SymReduction.orbit_select_minimalPeriod}{SYMLEAN/SymmetricReduction.lean:123}, \leanref{ErdosSzekeres.SymReduction.orbit_select}{SYMLEAN/SymmetricReduction.lean:194}]\label{sr:lem:orbit}
Let $X$ be a set and $f\colon X\to X$.
\begin{enumerate}[label=(\alph*)]
\item (\lean{orbit_select_minimalPeriod}) Let $p\in\N$ with $0<p$. For every finite $S\subseteq X$: if $f(y)\in S$ for
  every $y\in S$, $f$ is injective on $S$, and $\operatorname{minimalPeriod}(f,y)=p$ for every $y\in S$, then
  $p\mid\card S$, and for every $k\in\N$ with $kp\le\card S$ there is $T\subseteq S$ with $f(y)\in T$ for every $y\in
  T$ and $\card T=kp$.
\item (\lean{orbit_select}) Let $p$ be a prime. For every finite $S\subseteq X$: if $f(y)\in S$ for every $y\in S$,
  $f$ is injective on $S$, $f^{[p]}(y)=y$ for every $y\in S$, and $f(y)\ne y$ for every $y\in S$, then $p\mid\card S$,
  and for every $k\in\N$ with $kp\le\card S$ there is $T\subseteq S$ with $f(y)\in T$ for every $y\in T$ and $\card
  T=kp$.
\end{enumerate}
\end{lemma}

\begin{proof}
Both follow from Lemma~\ref{sr:lem:orbit-core}; we check its hypotheses. (a) For $y\in S$, $f^{[p]}(y)=y$ by (MP3) and
$\card{\operatorname{orbitFS}(f,p,y)}=p$ by Lemma~\ref{sr:lem:orbitFS}(e). (b) $p>0$ because $p$ is prime, and
$\card{\operatorname{orbitFS}(f,p,y)}=p$ for $y\in S$ by Lemma~\ref{sr:lem:orbitFS}(d).
\end{proof}

In (b) the injectivity hypothesis is redundant: if $f(y)=f(y')$ with $y,y'\in S$, then applying $f^{[p-1]}$ gives
$y=f^{[p]}(y)=f^{[p]}(y')=y'$. It is part of the Lean statement, and it is available in all applications.

\subsection{Rotations and reflections in coordinates}\label{sr:ssec:maps}

\begin{lemma}[The maps $\operatorname{crot}$; \leanref{ErdosSzekeres.SymReduction.toC_crot}{SYMLEAN/SymmetricReduction.lean:351}, \leanref{ErdosSzekeres.SymReduction.crot_one}{SYMLEAN/SymmetricReduction.lean:354}, \leanref{ErdosSzekeres.SymReduction.crot_comp}{SYMLEAN/SymmetricReduction.lean:359}, \leanref{ErdosSzekeres.SymReduction.crot_iterate}{SYMLEAN/SymmetricReduction.lean:364}, \leanref{ErdosSzekeres.SymReduction.crot_eq_self_iff}{SYMLEAN/SymmetricReduction.lean:370}, \leanref{ErdosSzekeres.SymReduction.crot_injective}{SYMLEAN/SymmetricReduction.lean:384}, \leanref{ErdosSzekeres.SymReduction.eDist2_crot_center}{SYMLEAN/SymmetricReduction.lean:394}]\label{sr:lem:crot}
Let $c,w,v\in\C$ and $p\in\R^2$.
\begin{enumerate}[label=(\alph*)]
\item (\lean{toC_crot}) $\operatorname{toC}(\operatorname{crot}_{c,w}(p))=c+w(\operatorname{toC}(p)-c)$.
\item (\lean{crot_one}) $\operatorname{crot}_{c,1}=\mathrm{id}$.
\item (\lean{crot_comp}) $\operatorname{crot}_{c,w}(\operatorname{crot}_{c,v}(p))=\operatorname{crot}_{c,wv}(p)$.
\item (\lean{crot_iterate}) $(\operatorname{crot}_{c,w})^{[k]}=\operatorname{crot}_{c,w^k}$ for every $k\in\N$.
\item (\lean{crot_eq_self_iff}) If $w\ne1$, then $\operatorname{crot}_{c,w}(p)=p$ if and only if
  $\operatorname{toC}(p)=c$.
\item (\lean{crot_injective}) If $w\ne0$, then $\operatorname{crot}_{c,w}$ is injective.
\item (\lean{eDist2_crot_center}) If $\mathrm{normSq}(w)=1$, then
  $\operatorname{eDist2}(\operatorname{crot}_{c,w}(p),\operatorname{ofC}(c))=
  \operatorname{eDist2}(p,\operatorname{ofC}(c))$.
\end{enumerate}
\end{lemma}

\begin{proof}
Write $z=\hat p$. Equalities of points are proved by applying the injective map $\operatorname{toC}$
(Lemma~\ref{sr:lem:complex}(c)). (a) is $\operatorname{toC}(\operatorname{ofC}(\zeta))=\zeta$. (b) $c+1\cdot(z-c)=z$.
(c) By (a), $c+w\bigl((c+v(z-c))-c\bigr)=c+wv(z-c)$. (d) Induction on $k$. For $k=0$ both sides are the identity, by
(b). If it holds for $k$, then $(\operatorname{crot}_{c,w})^{[k+1]}(p)=\operatorname{crot}_{c,w}(\operatorname{crot}_{c,w^k}(p))
=\operatorname{crot}_{c,w\cdot w^k}(p)=\operatorname{crot}_{c,w^{k+1}}(p)$ by (c). (e) By (a),
$\operatorname{toC}(\operatorname{crot}_{c,w}(p))-z=(w-1)(z-c)$. For $w\ne1$ this vanishes if and only if $z=c$. (f) If
$\operatorname{crot}_{c,w}(p)=\operatorname{crot}_{c,w}(q)$, then by (a) $w(\hat p-\hat q)=0$, so $\hat p=\hat q$ and
$p=q$. (g) By Lemma~\ref{sr:lem:complex}(c),(d) and (a), the left side is $\mathrm{normSq}(c+w(z-c)-c)=
\mathrm{normSq}(w)\,\mathrm{normSq}(z-c)=\mathrm{normSq}(z-c)$, which is the right side.
\end{proof}

\begin{lemma}[Rotations and the half-turn as $\operatorname{crot}$; \leanref{ErdosSzekeres.SymReduction.rotAbout_eq_crot}{SYMLEAN/SymmetricReduction.lean:404}, \leanref{ErdosSzekeres.SymReduction.halfTurn_eq_crot}{SYMLEAN/SymmetricReduction.lean:415}]\label{sr:lem:rot-half}
For $c\in\R^2$ and $\theta\in\R$:
\begin{enumerate}[label=(\alph*)]
\item (\lean{rotAbout_eq_crot}) $\operatorname{rotAbout}_{c,\theta}=\operatorname{crot}_{\operatorname{toC}(c),\,\exp(\theta i)}$ (equality of maps
  of $\R^2$);
\item (\lean{halfTurn_eq_crot}) $\operatorname{halfTurn}_c=\operatorname{crot}_{\operatorname{toC}(c),\,-1}$.
\end{enumerate}
\end{lemma}

\begin{proof}
Let $p\in\R^2$, $x=p_0-c_0$ and $y=p_1-c_1$, so that $\hat p-\hat c=x+iy$. (a) By Euler's formula
$\exp(\theta i)=\cos\theta+i\sin\theta$, hence $\exp(\theta i)(x+iy)=(x\cos\theta-y\sin\theta)+i(x\sin\theta+y\cos\theta)$.
Adding $\hat c=c_0+ic_1$ gives $\operatorname{toC}(\operatorname{rotAbout}_{c,\theta}(p))$; conclude by
Lemma~\ref{sr:lem:crot}(a) and the injectivity of $\operatorname{toC}$. (b) $\hat c+(-1)(\hat p-\hat c)=2\hat c-\hat p=
\operatorname{toC}(\operatorname{halfTurn}_c(p))$.
\end{proof}

\begin{lemma}[Line reflections; \leanref{ErdosSzekeres.SymReduction.toC_reflLine}{SYMLEAN/SymmetricReduction.lean:427}, \leanref{ErdosSzekeres.SymReduction.reflLine_involutive}{SYMLEAN/SymmetricReduction.lean:435}, \leanref{ErdosSzekeres.SymReduction.cross_of_reflLine_fixed}{SYMLEAN/SymmetricReduction.lean:447}, \leanref{ErdosSzekeres.SymReduction.orient_eq_zero_of_cross}{SYMLEAN/SymmetricReduction.lean:455}]\label{sr:lem:refl}
Let $a,u,p,q,r\in\R^2$.
\begin{enumerate}[label=(\alph*)]
\item (\lean{toC_reflLine}) If $u_0^2+u_1^2=1$, then $\operatorname{toC}(\operatorname{reflLine}_{a,u}(p))=\hat a+\hat u^{\,2}\,
  \overline{(\hat p-\hat a)}$.
\item (\lean{reflLine_involutive}) If $u_0^2+u_1^2=1$, then
  $\operatorname{reflLine}_{a,u}(\operatorname{reflLine}_{a,u}(p))=p$.
\item (\lean{cross_of_reflLine_fixed}) If $\operatorname{reflLine}_{a,u}(p)=p$, then $(p_0-a_0)u_1-(p_1-a_1)u_0=0$.
  (No hypothesis on $u$.)
\item (\lean{orient_eq_zero_of_cross}) If $u_0^2+u_1^2=1$ and $(z_0-a_0)u_1-(z_1-a_1)u_0=0$ for each
  $z\in\{p,q,r\}$, then $\orient(p,q,r)=0$.
\end{enumerate}
\end{lemma}

\begin{proof}
Let $x=p_0-a_0$, $y=p_1-a_1$ and $d=xu_0+yu_1$, so that $\operatorname{reflLine}_{a,u}(p)-a=(2du_0-x,\,2du_1-y)$.

(a) Expanding, $\hat u^{\,2}\,\overline{(x+iy)}=(u_0^2-u_1^2+2iu_0u_1)(x-iy)=\bigl((u_0^2-u_1^2)x+2u_0u_1y\bigr)+
i\bigl(2u_0u_1x-(u_0^2-u_1^2)y\bigr)$. On the other hand $2du_0-x=(2u_0^2-1)x+2u_0u_1y$ and
$2du_1-y=2u_0u_1x+(2u_1^2-1)y$. Since $u_0^2+u_1^2=1$, we have $2u_0^2-1=u_0^2-u_1^2$ and $2u_1^2-1=-(u_0^2-u_1^2)$, so the
two expressions agree. (In other words, the real and imaginary parts of the difference are $x(u_0^2+u_1^2-1)$ and
$y(u_0^2+u_1^2-1)$.)

(b) Put $\zeta=\hat p-\hat a$; note $\hat u\,\overline{\hat u}=u_0^2+u_1^2=1$. By (a), applied twice,
$\operatorname{toC}(\operatorname{reflLine}_{a,u}(\operatorname{reflLine}_{a,u}(p)))-\hat a=\hat u^{\,2}\,
\overline{\hat u^{\,2}\,\bar\zeta}=(\hat u\,\overline{\hat u})^2\zeta=\zeta$. Hence the two points have the same
$\operatorname{toC}$, and they are equal.

(c) The fixed-point equations read $a_0+2du_0-x=p_0=a_0+x$ and $a_1+2du_1-y=a_1+y$, that is, $du_0=x$ and $du_1=y$. Hence
$xu_1-yu_0=du_0u_1-du_1u_0=0$.

(d) Write $v\times w=v_0w_1-v_1w_0$ and $v\cdot w=v_0w_0+v_1w_1$ for $v,w\in\R^2$, so that
$\orient(p,q,r)=(q-p)\times(r-p)$. For all $v,w,u\in\R^2$ one has the identity
\[
  (v\times w)(u_0^2+u_1^2)=(w\cdot u)(v\times u)-(v\cdot u)(w\times u),
\]
checked by expanding both sides. Take $v=q-p$ and $w=r-p$. The hypotheses say $(z-a)\times u=0$ for $z=p,q,r$, so
$v\times u=(q-a)\times u-(p-a)\times u=0$, and likewise $w\times u=0$. Hence $\orient(p,q,r)\cdot1=0$.
\end{proof}

\subsection{Isometries of the plane}\label{sr:ssec:iso}

\begin{lemma}[Isometries in coordinates; \leanref{ErdosSzekeres.SymReduction.IsEucIsometry.injective}{SYMLEAN/SymmetricReduction.lean:301}, \leanref{ErdosSzekeres.SymReduction.isometry_toC_conj}{SYMLEAN/SymmetricReduction.lean:313}]\label{sr:lem:toC-iso}
Let $g\colon\R^2\to\R^2$ be a Euclidean isometry.
\begin{enumerate}[label=(\alph*)]
\item (\lean{IsEucIsometry.injective}) $g$ is injective.
\item (\lean{isometry_toC_conj}) The map $G\colon\C\to\C$, $G(z)=\operatorname{toC}(g(\operatorname{ofC}(z)))$, is
  an isometry of $\C$:
  $\abs{G(z)-G(w)}=\abs{z-w}$ for all $z,w\in\C$. (Lean: Mathlib's \texttt{Isometry G}; for the metric space $\C$, whose
  distance is $\abs{z-w}$, this is equivalent to the displayed equation.)
\end{enumerate}
\end{lemma}

\begin{proof}
(a) If $g(p)=g(q)$, then $\operatorname{eDist2}(p,q)=\operatorname{eDist2}(g(p),g(q))=0$, so $p=q$ by
Lemma~\ref{sr:lem:complex}(e). (b) By Lemma~\ref{sr:lem:complex}(c),(d),
\[
\abs{G(z)-G(w)}^2=\operatorname{eDist2}\bigl(g(\operatorname{ofC}z),g(\operatorname{ofC}w)\bigr)=
\operatorname{eDist2}(\operatorname{ofC}z,\operatorname{ofC}w)=\abs{z-w}^2,
\]
and both moduli are nonnegative.
\end{proof}

\begin{theorem}[Isometries of the plane; \leanref{ErdosSzekeres.SymReduction.isometry_complex_cases}{SYMLEAN/SymmetricReduction.lean:321}, \leanref{ErdosSzekeres.SymReduction.isometry_structure}{SYMLEAN/SymmetricReduction.lean:869}]\label{sr:thm:structure}
Part (a) describes the isometries of $\C$; part (b) the Euclidean isometries of the plane that map a finite set with at
least two points into itself.
\begin{enumerate}[label=(\alph*)]
\item (\lean{isometry_complex_cases}) Let $G\colon\C\to\C$ be an isometry: $\abs{G(z)-G(w)}=\abs{z-w}$ for all
  $z,w\in\C$. Then there are $a,b\in\C$ with $\abs a=1$ such that either $G(z)=az+b$ for all $z\in\C$, or
  $G(z)=a\bar z+b$ for all $z\in\C$.
\item (\lean{isometry_structure}) Let $S\subseteq\R^2$ be finite with $1<\card S$, and let $g\colon\R^2\to\R^2$ be a
  Euclidean isometry with $g\ne\mathrm{id}$ (as maps of $\R^2$) and $\operatorname{Invariant}(g,S)$. Then one of the
  following holds:
  \begin{enumerate}[label=(\roman*)]
  \item there are $c,a\in\C$ with $g=\operatorname{crot}_{c,a}$, $a\ne1$, and $a^m=1$ for some $m\in\N$ with $0<m$;
  \item there are $A,u\in\R^2$ with $u_0^2+u_1^2=1$ and $g=\operatorname{reflLine}_{A,u}$.
  \end{enumerate}
\end{enumerate}
\end{theorem}

\begin{proof}
\emph{(a).} This proof differs from the Lean proof. The Lean proof applies Mathlib's construction of an affine isometry
from an isometry between real normed affine spaces with strictly convex codomain
(\lean{Isometry.affineIsometryOfStrictConvexSpace}, a definition whose underlying map is the given isometry; it is a
result of Mazur--Ulam type that does not need surjectivity), and then Mathlib's theorem describing the real-linear
isometric equivalences of $\C$ (\lean{linear_isometry_complex}). We give a self-contained argument.

Let $b=G(0)$ and $H(z)=G(z)-b$. Then $H(0)=0$, $\abs{H(z)-H(w)}=\abs{z-w}$, and so $\abs{H(z)}=\abs{H(z)-H(0)}=\abs z$.
For all $z,w$ we have $\abs{z-w}^2=\abs z^2-2\langle z,w\rangle+\abs w^2$, and the same for $H(z),H(w)$. Hence
\[
2\langle H(z),H(w)\rangle=\abs{H(z)}^2+\abs{H(w)}^2-\abs{H(z)-H(w)}^2=\abs z^2+\abs w^2-\abs{z-w}^2=2\langle z,w\rangle,
\]
so $H$ preserves the inner product. Put $a=H(1)$ and $e=H(i)$. Then $\abs a=\abs1=1$, $\abs e=1$ and $\langle
e,a\rangle=\langle i,1\rangle=0$. Let $v=e\bar a$. Then $\abs v=1$ and $\operatorname{Re}v=\langle e,a\rangle=0$, so
$v=i$ or $v=-i$; and $e=e\bar a a=va$, so $e=ia$ or $e=-ia$.

Next, every $\xi\in\C$ satisfies $\xi=\langle\xi,a\rangle a+\langle\xi,e\rangle e$. Indeed, write $\bar a\xi=\alpha+i\beta$
with $\alpha,\beta\in\R$; then $\xi=a\bar a\xi=\alpha a+\beta(ia)$, $\langle\xi,a\rangle=\operatorname{Re}(\xi\bar a)=\alpha$
and $\langle\xi,ia\rangle=\operatorname{Re}(-i\xi\bar a)=\beta$; and $\langle\xi,e\rangle e=\langle\xi,ia\rangle\,ia$ for
both choices $e=\pm ia$.

Now let $z=x+iy$ with $x,y\in\R$. Since $H$ preserves the inner product, $\langle H(z),a\rangle=\langle
H(z),H(1)\rangle=\langle z,1\rangle=x$ and $\langle H(z),e\rangle=\langle z,i\rangle=y$. Hence $H(z)=xa+ye$. If $e=ia$,
this is $a(x+iy)=az$; if $e=-ia$, it is $a(x-iy)=a\bar z$. The choice between the two cases does not depend on $z$. So
$G(z)=az+b$ for all $z$, or $G(z)=a\bar z+b$ for all $z$, and $\abs a=1$.

\emph{(b).} This part follows the Lean proof. Let $G(z)=\operatorname{toC}(g(\operatorname{ofC}z))$. By
Lemma~\ref{sr:lem:toC-iso}(b) and part (a) there are $a,b\in\C$ with $\abs a=1$ such that either
\textup{(R)} $G(z)=az+b$ for all $z$, or \textup{(C)} $G(z)=a\bar z+b$ for all $z$. Since $\operatorname{ofC}(\hat
p)=p$, we have $\operatorname{toC}(g(p))=G(\hat p)$ for every $p\in\R^2$. The set $S$ is nonempty because $\card S>1$.

\emph{Case (R), $a=1$.} Then $\operatorname{toC}(g(p))=\hat p+b$ for all $p$, and by induction
$\operatorname{toC}(g^{[k]}(q))=\hat q+kb$ for all $k\in\N$ and $q\in\R^2$. Take $q\in S$. By Lemma~\ref{sr:lem:pigeon}
there are $i<j$ with $g^{[i]}(q)=g^{[j]}(q)$, so $\hat q+ib=\hat q+jb$ and $(j-i)b=0$. Since $j-i$ is a nonzero
integer, $b=0$. Then $\operatorname{toC}(g(p))=\hat p$ for all $p$, so $g=\mathrm{id}$ by the injectivity of
$\operatorname{toC}$. This contradicts $g\ne\mathrm{id}$; so this case does not occur.

\emph{Case (R), $a\ne1$.} Put $c=b/(1-a)$. For every $p$, $c+a(\hat p-c)=a\hat p+(1-a)c=a\hat p+b=
\operatorname{toC}(g(p))$, so $g=\operatorname{crot}_{c,a}$ by Lemma~\ref{sr:lem:crot}(a). Also $a\ne0$ because
$\abs a=1$. Since $\card S\ge2$, some $q\in S$ differs from $\operatorname{ofC}(c)$. By Lemma~\ref{sr:lem:pigeon},
applied to $g=\operatorname{crot}_{c,a}$, there are $i<j$ with $(\operatorname{crot}_{c,a})^{[i]}(q)=
(\operatorname{crot}_{c,a})^{[j]}(q)$. By Lemma~\ref{sr:lem:crot}(d),(a), $c+a^i(\hat q-c)=c+a^j(\hat q-c)$. Here
$\hat q-c\ne0$, since otherwise $q=\operatorname{ofC}(\hat q)=\operatorname{ofC}(c)$. Hence $a^i=a^j=a^ia^{j-i}$, and as
$a^i\ne0$, $a^{j-i}=1$. With $m=j-i>0$ we obtain alternative (i).

\emph{Case (C).} We prove alternative (ii); the hypothesis $g\ne\mathrm{id}$ is not needed here. Since $\abs a=1$,
$a\bar a=1$. For every $p$,
\[
\operatorname{toC}(g(g(p)))=a\,\overline{(a\bar{\hat p}+b)}+b=a\bar a\hat p+a\bar b+b=\hat p+t,\qquad t:=a\bar b+b .
\]
By induction, $\operatorname{toC}((g\circ g)^{[k]}(p))=\hat p+kt$ for all $k$ and $p$. The map $g\circ g$ maps $S$ into
$S$. Take $q\in S$; Lemma~\ref{sr:lem:pigeon} applied to $g\circ g$ gives $i<j$ with $(j-i)t=0$, hence $t=0$, that is,
$a\bar b=-b$. Choose $u\in\C$ with $u^2=a$ (every complex number has a square root; the Lean proof takes the principal
value $a^{1/2}$). Then $\abs u^2=\abs a=1$, so the point $U=\operatorname{ofC}(u)$ satisfies $U_0^2+U_1^2=1$. Let
$A=\operatorname{ofC}(b/2)$. By Lemma~\ref{sr:lem:refl}(a), for every $p$,
\[
\operatorname{toC}(\operatorname{reflLine}_{A,U}(p))=\tfrac b2+u^2\,\overline{\bigl(\hat p-\tfrac b2\bigr)}
=a\bar{\hat p}+\tfrac b2-\tfrac{a\bar b}{2}=a\bar{\hat p}+\tfrac b2+\tfrac b2=a\bar{\hat p}+b=\operatorname{toC}(g(p)).
\]
Hence $g=\operatorname{reflLine}_{A,U}$, which is alternative (ii).
\end{proof}

In (b), the pigeonhole steps are where finiteness enters: they exclude translations ($a=1$) and glide reflections
($t\ne0$), and they make $a$ a root of unity.

\subsection{Elementary geometric lemmas}\label{sr:ssec:geom}

\begin{lemma}[Strict supporting lines; \leanref{ErdosSzekeres.Suk.convexPosition_of_strict_supports}{ErdosSzekeres/Suk/SupportingLines.lean:28}]\label{sr:lem:supports}
Let $s\subseteq\R^2$ be finite. If for every $p\in s$ there is $q\in\R^2$ with $\orient(p,q,r)>0$ for every $r\in s$
with $r\ne p$, then $s$ is in convex position.
\end{lemma}

\begin{proof}
This is the lemma titled ``Strict supporting lines'' in the foundations section, the same declaration of the main
library; the short proof is repeated here because Lemma~\ref{sr:lem:concyclic} uses it. By Lemma~\ref{fd:lem:cp-hull} it
suffices to show $p\notin\conv(s\setminus\{p\})$ for every $p\in s$. Fix $p$ and a point $q$ as in the hypothesis, and
let $\varphi(x)=\orient(p,q,x)=(q_0-p_0)(x_1-p_1)-(q_1-p_1)(x_0-p_0)$. This is an affine function of $x$:
$\varphi(\lambda x+\mu y)=\lambda\varphi(x)+\mu\varphi(y)$ whenever $\lambda+\mu=1$. So the open half-plane
$H=\{x:\varphi(x)>0\}$ is convex: for $x,y\in H$ and $\lambda,\mu\ge0$ with $\lambda+\mu=1$, both terms
$\lambda\varphi(x)$ and $\mu\varphi(y)$ are $\ge0$ and at least one is $>0$. By hypothesis $s\setminus\{p\}\subseteq
H$, hence $\conv(s\setminus\{p\})\subseteq H$. But $\varphi(p)=\orient(p,q,p)=0$, so $p\notin H$, and therefore
$p\notin\conv(s\setminus\{p\})$.
\end{proof}

\begin{lemma}[Concyclic points; \leanref{ErdosSzekeres.SymReduction.convexPosition_of_concyclic}{SYMLEAN/SymmetricReduction.lean:509}]\label{sr:lem:concyclic}
Let $s\subseteq\R^2$ be finite, $c\in\R^2$ and $r\in\R$, and assume $\operatorname{eDist2}(p,c)=r$ for every $p\in s$.
Then $s$ is in convex position.
\end{lemma}

No general position is assumed, and $r$ is arbitrary (for $r=0$ the set has at most one point).

\begin{proof}
We apply Lemma~\ref{sr:lem:supports} with the tangent direction. For $p\in s$ let
$q=\bigl(p_0-(p_1-c_1),\;p_1+(p_0-c_0)\bigr)$, that is, $p$ plus the vector $p-c$ turned by $+90^\circ$. For every
$x\in\R^2$, expanding gives
\[
\orient(p,q,x)=-\bigl[(p_0-c_0)(x_0-p_0)+(p_1-c_1)(x_1-p_1)\bigr]
=\tfrac12\bigl(\operatorname{eDist2}(x,p)+\operatorname{eDist2}(p,c)-\operatorname{eDist2}(x,c)\bigr).
\]
(The second equality is $\abs{x-c}^2=\abs{x-p}^2+2\langle x-p,p-c\rangle+\abs{p-c}^2$.) For $x\in s$ we have
$\operatorname{eDist2}(x,c)=r=\operatorname{eDist2}(p,c)$, so $\orient(p,q,x)=\frac12\operatorname{eDist2}(x,p)$, which is
positive when $x\ne p$ (Lemma~\ref{sr:lem:complex}(e)). Lemma~\ref{sr:lem:supports} gives the claim.
\end{proof}

\begin{corollary}[A concyclic subset gives convex polygons; \leanref{ErdosSzekeres.SymReduction.hasConvexNGon_of_concyclic}{SYMLEAN/SymmetricReduction.lean:522}]\label{sr:cor:concyclic}
Let $T\subseteq S$ be finite sets of points, $c\in\R^2$ and $r\in\R$ with $\operatorname{eDist2}(p,c)=r$ for every $p\in
T$, and let $n\in\N$ with $n\le\card T$. Then \texttt{HasConvexNGon n S}.
\end{corollary}

\begin{proof}
$T$ is in convex position by Lemma~\ref{sr:lem:concyclic}, hence \texttt{HasConvexNGon n T} since $n\le\card T$, hence
\texttt{HasConvexNGon n S} by monotonicity (Lemma~\ref{fd:lem:mono}).
\end{proof}

\begin{lemma}[The centre of a half-turn; \leanref{ErdosSzekeres.SymReduction.center_not_mem_of_halfTurn}{SYMLEAN/SymmetricReduction.lean:528}]\label{sr:lem:center}
Let $S$ be in general position with $2\le\card S$, let $c\in\R^2$, and assume
$\operatorname{Invariant}(\operatorname{halfTurn}_c,S)$. Then $c\notin S$.
\end{lemma}

\begin{proof}
Suppose $c\in S$. Since $\card S\ge2$, there is $q\in S$ with $q\ne c$, and $q'=\operatorname{halfTurn}_c(q)=2c-q$ lies
in $S$. We have $q'\ne q$, since $2c-q=q$ would give $q=c$, and $q'\ne c$, since $2c-q=c$ would give $q=c$. So $q,c,q'$
are pairwise distinct points of $S$. But
\[
\orient(q,c,2c-q)=(c_0-q_0)\cdot2(c_1-q_1)-(c_1-q_1)\cdot2(c_0-q_0)=0,
\]
contradicting general position (Lemma~\ref{fd:lem:collinear}).
\end{proof}

\subsection{Roots of unity}\label{sr:ssec:roots}

\begin{definition}[Primitive roots of unity, order; written step]\label{sr:def:prim}
For $\zeta\in\C$ and $k\in\N$, $\zeta$ is a \emph{primitive $k$-th root of unity} (Mathlib \lean{IsPrimitiveRoot}) if
$\zeta^k=1$ and every $l\in\N$ with $\zeta^l=1$ is a multiple of $k$. If $\zeta^m=1$ for some $m>0$, the \emph{order}
$\operatorname{ord}(\zeta)$ (Mathlib \lean{orderOf}) is the least such $m$; otherwise it is $0$.
\end{definition}

\begin{lemma}[Facts on roots of unity; written step]\label{sr:lem:prim}
Let $\zeta,\xi\in\C$ and $n,k,d\in\N$.
\begin{enumerate}[label=(P\arabic*)]
\item If $\zeta^m=1$ for some $m>0$, then $\operatorname{ord}(\zeta)>0$, $\zeta$ is a primitive
  $\operatorname{ord}(\zeta)$-th root of unity, and $\operatorname{ord}(\zeta)=1$ if and only if $\zeta=1$. (Mathlib
  \lean{IsPrimitiveRoot.orderOf}, \lean{orderOf_eq_one_iff}.)
\item If $\zeta$ is a primitive $n$-th root of unity, $n>0$ and $n=k\,d$, then $\zeta^k$ is a primitive $d$-th root of
  unity. (\lean{IsPrimitiveRoot.pow}.)
\item A primitive $2$nd root of unity equals $-1$. (\lean{IsPrimitiveRoot.eq_neg_one_of_two_right}.)
\item If $\zeta$ is a primitive $k$-th root of unity and $k>1$, then $\zeta^k=1$, $\zeta\ne1$ and $\zeta\ne0$.
  (\lean{IsPrimitiveRoot.pow_eq_one}, \lean{IsPrimitiveRoot.ne_one}, \lean{IsPrimitiveRoot.ne_zero}.)
\item If $\zeta$ is a primitive $k$-th root of unity, $k>0$ and $\xi^k=1$, then $\xi=\zeta^e$ for some $e<k$.
  (\lean{IsPrimitiveRoot.eq_pow_of_pow_eq_one}.)
\item If $n>0$, then $\exp(2\pi i/n)$ is a primitive $n$-th root of unity. (\lean{Complex.isPrimitiveRoot_exp}.)
\end{enumerate}
\end{lemma}

\begin{proof}
(P1) Put $s=\operatorname{ord}(\zeta)$. Since $\zeta^m=1$ with $m>0$, Definition~\ref{sr:def:prim} makes $s$ the least
positive integer with $\zeta^s=1$; in particular $s>0$. If $\zeta^l=1$, write $l=ts+r$ with $0\le r<s$; then
$1=\zeta^l=(\zeta^s)^t\zeta^r=\zeta^r$, so $r=0$ by minimality, and $s\mid l$. So $\zeta$ is a primitive $s$-th root of
unity. Finally $s=1$ if and only if $\zeta^1=1$, that is, $\zeta=1$.

(P2) $(\zeta^k)^d=\zeta^n=1$. If $(\zeta^k)^l=1$, then $\zeta^{kl}=1$, so $kd=n\mid kl$; since $n>0$ we have $k>0$, and
so $d\mid l$.

(P3) Let $\zeta$ be a primitive $2$nd root. Then $\zeta^2=1$, and $\zeta\ne1$ because $\zeta^1=1$ would give $2\mid1$.
From $(\zeta-1)(\zeta+1)=\zeta^2-1=0$ we get $\zeta=-1$.

(P4) $\zeta^k=1$ by definition. If $\zeta=1$, then $\zeta^1=1$ gives $k\mid1$, so $k=1$. If $\zeta=0$, then
$\zeta^k=0\ne1$ since $k\ge1$.

(P5) The numbers $\zeta^0,\dots,\zeta^{k-1}$ are pairwise distinct: if $\zeta^i=\zeta^j$ with $i<j<k$, then
$\zeta^i\ne0$ (as $\zeta^k=1$), so $\zeta^{j-i}=1$ and $k\mid j-i$ with $0<j-i<k$, which is impossible. They are $k$
distinct roots of the polynomial $X^k-1$, which has degree $k$ and therefore at most $k$ roots in $\C$. Since $\xi$ is
also a root, it is one of them.

(P6) Let $\zeta=\exp(2\pi i/n)$. For $l\in\N$, $\zeta^l=\exp(2\pi il/n)$, and $\exp(w)=1$ holds exactly when $w\in2\pi
i\Z$. So $\zeta^l=1$ if and only if $l/n\in\Z$, that is, $n\mid l$. In particular $\zeta^n=1$.
\end{proof}

\begin{lemma}[The angles $2\pi/3$, $2\pi/4$, $2\pi/5$; \leanref{ErdosSzekeres.SymReduction.isPrimitiveRoot_rot3}{SYMLEAN/SymmetricReduction.lean:584}, \leanref{ErdosSzekeres.SymReduction.isPrimitiveRoot_rot5}{SYMLEAN/SymmetricReduction.lean:591}, \leanref{ErdosSzekeres.SymReduction.isPrimitiveRoot_rot4}{SYMLEAN/SymmetricReduction.lean:598}]\label{sr:lem:rot345}
For $n\in\{3,5,4\}$, $\exp\bigl((2\pi/n)\,i\bigr)$ is a primitive $n$-th root of unity. Here $2\pi/n$ is the real
number $2\pi/n$ (Lean \texttt{2 * Real.pi / n}), viewed as a complex number.
\end{lemma}

\begin{proof}
$(2\pi/n)i=2\pi i/n$; apply (P6).
\end{proof}

\begin{lemma}[Even order or an odd prime; \leanref{ErdosSzekeres.SymReduction.root_cases}{SYMLEAN/SymmetricReduction.lean:552}]\label{sr:lem:root-cases}
Let $a\in\C$ and $m\in\N$ with $0<m$, $a^m=1$ and $a\ne1$. Then either $a^k=-1$ for some $k\in\N$, or there are
$p,k\in\N$ such that $p$ is prime, $p\ne2$, and $a^k$ is a primitive $p$-th root of unity.
\end{lemma}

\begin{proof}
Let $n=\operatorname{ord}(a)$. By (P1), $n>0$, $a$ is a primitive $n$-th root of unity, and $n\ne1$ because $a\ne1$.
If $2\mid n$, write $n=2k=k\cdot2$; by (P2), $a^k$ is a primitive $2$nd root of unity, so $a^k=-1$ by (P3). Otherwise
let $p$ be the least prime factor of $n$ (it exists since $n>1$). Then $p\ne2$ because $n$ is odd, and $n=(n/p)\,p$, so
$a^{n/p}$ is a primitive $p$-th root of unity by (P2).
\end{proof}

In the proof the first alternative is taken for \emph{every} even order (so orders $6,10,12,\dots$ lead to the
half-turn), and in the second $p$ is the least prime factor of the order; the statement itself does not record this.

\begin{lemma}[Odd primes; \leanref{ErdosSzekeres.SymReduction.prime_cases}{SYMLEAN/SymmetricReduction.lean:572}]\label{sr:lem:prime-cases}
If $p$ is prime and $p\ne2$, then $p=3$, $p=5$ or $7\le p$.
\end{lemma}

\begin{proof}
$p\ge2$ and $p\ne2$, so $p\in\{3,4,5,6\}$ or $p\ge7$; and $4=2\cdot2$ and $6=2\cdot3$ are not prime.
\end{proof}

\begin{lemma}[Order $2$, a multiple of $4$, or an odd prime; \leanref{ErdosSzekeres.SymReduction.root_cases2}{SYMLEAN/SymmetricReduction.lean:607}]\label{sr:lem:root-cases2}
Let $a\in\C$ and $m\in\N$ with $0<m$, $a^m=1$ and $a\ne1$. Then $a=-1$, or $a^k$ is a primitive $4$th root of unity for
some $k\in\N$, or there are $p,k\in\N$ such that $p$ is prime, $p\ne2$, and $a^k$ is a primitive $p$-th root of unity.
\end{lemma}

\begin{proof}
Let $n=\operatorname{ord}(a)$; as before, $n>1$ and $a$ is a primitive $n$-th root of unity. If some prime $p\ne2$
divides $n$, then $a^{n/p}$ is a primitive $p$-th root of unity by (P2). Otherwise $2$ is the only prime factor of $n$.
If $4\mid n$, then $a^{n/4}$ is a primitive $4$th root of unity by (P2). Otherwise $n$ is even (its least prime factor
is $2$), say $n=2m'$ with $m'\ge1$. If $m'>1$, a prime factor of $m'$ divides $n$, so it equals $2$, and then $4\mid
n$, which is excluded. So $m'=1$, $n=2$, and $a=-1$ by (P3).
\end{proof}

\subsection{The case lemmas}\label{sr:ssec:cases}

\begin{lemma}[Half-turn invariant sets; \leanref{ErdosSzekeres.SymReduction.halfTurn_facts}{SYMLEAN/SymmetricReduction.lean:647}]\label{sr:lem:halfturn-facts}
Let $S$ be in general position with $2\le\card S$, let $c\in\R^2$, and assume
$\operatorname{Invariant}(\operatorname{halfTurn}_c,S)$. Then $c\notin S$, $2\mid\card S$, and for every $k\in\N$ with
$2k\le\card S$ there is $T\subseteq S$ with $\operatorname{Invariant}(\operatorname{halfTurn}_c,T)$ and $\card T=2k$.
\end{lemma}

\begin{proof}
$c\notin S$ by Lemma~\ref{sr:lem:center}. Let $f=\operatorname{halfTurn}_c=\operatorname{crot}_{\hat c,-1}$
(Lemma~\ref{sr:lem:rot-half}(b)). By Lemma~\ref{sr:lem:crot}(d),(b), $f^{[2]}=\operatorname{crot}_{\hat c,1}=\mathrm{id}$.
If $f(y)=y$ for some $y\in S$, then $\hat y=\hat c$ by Lemma~\ref{sr:lem:crot}(e) (as $-1\ne1$), so $y=c\in S$, a
contradiction. $f$ is injective by Lemma~\ref{sr:lem:crot}(f). Lemma~\ref{sr:lem:orbit}(b) with the prime $p=2$ gives
the claims.
\end{proof}

\begin{lemma}[Case of a half-turn; \leanref{ErdosSzekeres.SymReduction.case_halfTurn}{SYMLEAN/SymmetricReduction.lean:670}; conditional on \lean{HalfTurn34}]\label{sr:lem:case-halfturn}
Assume \lean{HalfTurn34} (a hypothesis). Let $S$ be in general position with $33\le\card S$, let $c\in\R^2$, and assume
$\operatorname{Invariant}(\operatorname{halfTurn}_c,S)$. Then \texttt{HasConvexNGon 7 S}.
\end{lemma}

\begin{proof}
By Lemma~\ref{sr:lem:halfturn-facts}, $c\notin S$ and $\card S$ is even; with $\card S\ge33$ this gives $\card S\ge34=17\cdot2$.
So there is $T\subseteq S$, invariant under $\operatorname{halfTurn}_c$, with $\card T=34$. $T$ is in general position
(monotonicity) and $c\notin T$. \lean{HalfTurn34} applied to $T$ and $c$ gives \texttt{HasConvexNGon 7 T}, hence
\texttt{HasConvexNGon 7 S} by monotonicity.
\end{proof}

\begin{lemma}[Rotations of prime order off the centre; \leanref{ErdosSzekeres.SymReduction.rot_erase_facts}{SYMLEAN/SymmetricReduction.lean:676}]\label{sr:lem:rot-erase}
Let $S$ be finite, $c,w\in\C$, and $p$ a prime with $w^p=1$ and $w\ne1$; assume
$\operatorname{Invariant}(\operatorname{crot}_{c,w},S)$. Let $S'=S\setminus\{\operatorname{ofC}(c)\}$. Then
$p\mid\card{S'}$, and for every $k\in\N$ with $kp\le\card{S'}$ there is $T\subseteq S$ with
$\operatorname{ofC}(c)\notin T$, $\operatorname{Invariant}(\operatorname{crot}_{c,w},T)$ and $\card T=kp$.
\end{lemma}

\begin{proof}
Let $f=\operatorname{crot}_{c,w}$. Since $p>0$ and $w^p=1$, $w\ne0$, so $f$ is injective (Lemma~\ref{sr:lem:crot}(f)).
The centre is fixed: $\operatorname{toC}(f(\operatorname{ofC}c))=c+w(c-c)=c$, so $f(\operatorname{ofC}c)=\operatorname{ofC}c$.
\emph{$S'$ is mapped into itself:} for $y\in S'$, $f(y)\in S$, and $f(y)=\operatorname{ofC}(c)=f(\operatorname{ofC}c)$
would give $y=\operatorname{ofC}(c)$ by injectivity. Further $f^{[p]}=\operatorname{crot}_{c,w^p}=\operatorname{crot}_{c,1}=
\mathrm{id}$ (Lemma~\ref{sr:lem:crot}(d),(b)). If $f(y)=y$ for some $y\in S'$, then $\hat y=c$ by
Lemma~\ref{sr:lem:crot}(e), so $y=\operatorname{ofC}(c)\notin S'$, a contradiction. Lemma~\ref{sr:lem:orbit}(b) applied
to $S'$ gives $p\mid\card{S'}$ and, for $kp\le\card{S'}$, an $f$-invariant $T\subseteq S'$ with $\card T=kp$; then
$T\subseteq S$ and $\operatorname{ofC}(c)\notin T$.
\end{proof}

\begin{lemma}[Rotations of prime order at least $7$; \leanref{ErdosSzekeres.SymReduction.case_rot_big}{SYMLEAN/SymmetricReduction.lean:711}]\label{sr:lem:rot-big}
Let $S$ be finite, $c,w\in\C$, and $p$ a prime with $7\le p$, $w^p=1$ and $w\ne1$. Assume
$\operatorname{Invariant}(\operatorname{crot}_{c,w},S)$ and $2\le\card S$. Then \texttt{HasConvexNGon 7 S}.
\end{lemma}

No general position is assumed.

\begin{proof}
Let $f=\operatorname{crot}_{c,w}$. Since $\card S\ge2$, some $q\in S$ differs from $\operatorname{ofC}(c)$. We have
$\mathrm{normSq}(w)^p=\mathrm{normSq}(w^p)=1$ and $\mathrm{normSq}(w)\ge0$; as $t\mapsto t^p$ is strictly increasing on
$[0,\infty)$, $\mathrm{normSq}(w)=1$. Let $O=\operatorname{orbitFS}(f,p,q)$. Then $O\subseteq S$
(Lemma~\ref{sr:lem:orbitFS}(b)); $f^{[p]}(q)=\operatorname{crot}_{c,w^p}(q)=q$; and $f(q)\ne q$, since otherwise $\hat
q=c$ by Lemma~\ref{sr:lem:crot}(e) and $q=\operatorname{ofC}(c)$. By Lemma~\ref{sr:lem:orbitFS}(d), $\card O=p$. Every
$x\in O$ has the form $x=f^{[i]}(q)=\operatorname{crot}_{c,w^i}(q)$ with $\mathrm{normSq}(w^i)=\mathrm{normSq}(w)^i=1$, so
$\operatorname{eDist2}(x,\operatorname{ofC}c)=\operatorname{eDist2}(q,\operatorname{ofC}c)$ by
Lemma~\ref{sr:lem:crot}(g). Corollary~\ref{sr:cor:concyclic} with $T=O$, centre $\operatorname{ofC}(c)$, radius
$\operatorname{eDist2}(q,\operatorname{ofC}c)$ and $n=7\le p=\card O$ gives \texttt{HasConvexNGon 7 S}.
\end{proof}

\begin{lemma}[Case of an odd prime order; \leanref{ErdosSzekeres.SymReduction.case_odd_prime}{SYMLEAN/SymmetricReduction.lean:736}; conditional on \lean{Threefold33} and \lean{Fivefold35}]\label{sr:lem:odd-prime}
Assume \lean{Threefold33} and \lean{Fivefold35} (hypotheses). Let $S$ be in general position with $33\le\card S$, let
$c,w\in\C$, let $p$ be a prime with $p\ne2$, let $w$ be a primitive $p$-th root of unity, and assume
$\operatorname{Invariant}(\operatorname{crot}_{c,w},S)$. Then \texttt{HasConvexNGon 7 S}.
\end{lemma}

\begin{proof}
By (P4), $w^p=1$ and $w\ne1$. Let $S'=S\setminus\{\operatorname{ofC}(c)\}$; then $\card{S'}\ge\card S-1\ge32$. By
Lemma~\ref{sr:lem:prime-cases}, $p=3$, $p=5$ or $p\ge7$.

\emph{$p=3$.} By Lemma~\ref{sr:lem:rot-erase}, $3\mid\card{S'}$; with $\card{S'}\ge32$ this gives $\card{S'}\ge33=11\cdot3$.
So there is $T\subseteq S$ with $\operatorname{ofC}(c)\notin T$, $\operatorname{Invariant}(\operatorname{crot}_{c,w},T)$
and $\card T=33$. Let $\zeta=\exp((2\pi/3)i)$. By Lemma~\ref{sr:lem:rot345}, $\zeta^3=1$, so by (P5) $\zeta=w^e$ for
some $e$. By Lemmas~\ref{sr:lem:rot-half}(a), \ref{sr:lem:complex}(c) and~\ref{sr:lem:crot}(d),
\[
\operatorname{rotAbout}_{\operatorname{ofC}(c),2\pi/3}=\operatorname{crot}_{\operatorname{toC}(\operatorname{ofC}c),\,\zeta}
=\operatorname{crot}_{c,w^e}=(\operatorname{crot}_{c,w})^{[e]},
\]
so $\operatorname{Invariant}(\operatorname{rotAbout}_{\operatorname{ofC}(c),2\pi/3},T)$ by Lemma~\ref{sr:lem:iterate}(b).
$T$ is in general position (monotonicity), $\card T=33$ and $\operatorname{ofC}(c)\notin T$. \lean{Threefold33}, applied
to $T$ and the centre $\operatorname{ofC}(c)$, gives \texttt{HasConvexNGon 7 T}, hence \texttt{HasConvexNGon 7 S}.

\emph{$p=5$.} In the same way, $5\mid\card{S'}$ and $\card{S'}\ge32$ give $\card{S'}\ge35=7\cdot5$; take $T\subseteq S$
with $\operatorname{ofC}(c)\notin T$, invariant under $\operatorname{crot}_{c,w}$, with $\card T=35$; with
$\zeta=\exp((2\pi/5)i)=w^e$, $T$ is invariant under $\operatorname{rotAbout}_{\operatorname{ofC}(c),2\pi/5}$; and
\lean{Fivefold35} gives \texttt{HasConvexNGon 7 T}, hence \texttt{HasConvexNGon 7 S}.

\emph{$p\ge7$.} Lemma~\ref{sr:lem:rot-big} applies, since $\card S\ge2$.
\end{proof}

\begin{lemma}[Case of a primitive fourth root; \leanref{ErdosSzekeres.SymReduction.case_fourfold}{SYMLEAN/SymmetricReduction.lean:763}; conditional on \lean{Fourfold32}]\label{sr:lem:fourfold}
Assume \lean{Fourfold32} (a hypothesis). Let $S$ be in general position with $33\le\card S$, let $c,w\in\C$, let $w$ be
a primitive $4$th root of unity, and assume $\operatorname{Invariant}(\operatorname{crot}_{c,w},S)$. Then
\texttt{HasConvexNGon 7 S}.
\end{lemma}

\begin{proof}
By (P4), $w^4=1$, $w\ne1$ and $w\ne0$; by (P2) with $4=2\cdot2$ and (P3), $w^2=-1$. Let $f=\operatorname{crot}_{c,w}$.
By Lemmas~\ref{sr:lem:iterate}(b) and~\ref{sr:lem:crot}(d), $S$ is invariant under $f^{[2]}=\operatorname{crot}_{c,-1}$,
which equals $\operatorname{halfTurn}_{\operatorname{ofC}(c)}$ by Lemma~\ref{sr:lem:rot-half}(b) and
$\operatorname{toC}(\operatorname{ofC}c)=c$. By Lemma~\ref{sr:lem:center}, $\operatorname{ofC}(c)\notin S$. Hence $\hat
y\ne c$ for every $y\in S$.

\emph{Every $y\in S$ has minimal period $4$ under $f$.} First, $f^{[4]}(y)=\operatorname{crot}_{c,w^4}(y)=y$, so by (MP1)
$d=\operatorname{minimalPeriod}(f,y)$ divides $4$, and $d\in\{1,2,4\}$. If $d=1$, then $f(y)=y$ by (MP2), and $\hat y=c$
by Lemma~\ref{sr:lem:crot}(e), which is impossible. If $d=2$, then $f^{[2]}(y)=y$ by (MP3), that is,
$\operatorname{crot}_{c,-1}(y)=y$, and again $\hat y=c$ (as $-1\ne1$). So $d=4$.

$f$ is injective (Lemma~\ref{sr:lem:crot}(f)). Lemma~\ref{sr:lem:orbit}(a) with $p=4$ gives $4\mid\card S$ and invariant
subsets of every size $4k\le\card S$. Since $8\cdot4=32\le\card S$, there is $T\subseteq S$ with
$\operatorname{Invariant}(f,T)$ and $\card T=32$. (Also $\card S\ge36$, but this is not needed.) Let
$\zeta=\exp((2\pi/4)i)$; by Lemma~\ref{sr:lem:rot345} and (P5), $\zeta=w^e$ for some $e$, and as in the proof of
Lemma~\ref{sr:lem:odd-prime}, $\operatorname{rotAbout}_{\operatorname{ofC}(c),2\pi/4}=f^{[e]}$, so $T$ is invariant under
it. $T$ is in general position, $\card T=32$, and $\operatorname{ofC}(c)\notin T$ because $T\subseteq S$.
\lean{Fourfold32} gives \texttt{HasConvexNGon 7 T}, hence \texttt{HasConvexNGon 7 S}.
\end{proof}

\begin{lemma}[Case of a reflection; \leanref{ErdosSzekeres.SymReduction.case_mirror}{SYMLEAN/SymmetricReduction.lean:811}; conditional on the two mirror classes]\label{sr:lem:mirror}
Assume \lean{MirrorOneAxis33} and \lean{MirrorNoAxis34} (hypotheses). Let $a,u\in\R^2$ with $u_0^2+u_1^2=1$, let $S$ be in
general position with $33\le\card S$, and assume $\operatorname{Invariant}(\operatorname{reflLine}_{a,u},S)$. Then
\texttt{HasConvexNGon 7 S}.
\end{lemma}

\begin{proof}
Let $\sigma=\operatorname{reflLine}_{a,u}$, $F=\{p\in S:\sigma(p)=p\}$ and $M=\{p\in S:\sigma(p)\ne p\}$, so that
$\card F+\card M=\card S$.

\emph{$\card F\le2$.} Otherwise $F$ contains three pairwise distinct points $x,y,z$. By Lemma~\ref{sr:lem:refl}(c),
each $t\in\{x,y,z\}$ satisfies $(t_0-a_0)u_1-(t_1-a_1)u_0=0$, so $\orient(x,y,z)=0$ by Lemma~\ref{sr:lem:refl}(d),
contradicting general position.

\emph{$M$ is mapped into itself.} By Lemma~\ref{sr:lem:refl}(b), $\sigma\circ\sigma=\mathrm{id}$; in particular $\sigma$
is injective. Let $y\in M$. Then $\sigma(y)\in S$, and if $\sigma(\sigma(y))=\sigma(y)$ then $y=\sigma(y)$, which is
false; so $\sigma(y)\in M$. Moreover $\sigma^{[2]}(y)=y$ and $\sigma(y)\ne y$ for $y\in M$. Lemma~\ref{sr:lem:orbit}(b)
with the prime $2$ gives: $2\mid\card M$, and for every $k$ with $2k\le\card M$ there is $T\subseteq M$, mapped into
itself by $\sigma$, with $\card T=2k$.

\emph{Case $\card F=0$.} Then $\card M=\card S\ge33$ and $\card M$ is even, so $\card M\ge34$. Take $T\subseteq M$
invariant with $\card T=34$. $T\subseteq S$ is in general position, and no point of $T$ is fixed by $\sigma$ (as
$T\subseteq M$). \lean{MirrorNoAxis34} applied to $T$, $a$, $u$ gives \texttt{HasConvexNGon 7 T}, hence
\texttt{HasConvexNGon 7 S}.

\emph{Case $\card F\ge1$.} Pick $x\in F$. Since $\card F\le2$, $\card M\ge31$, and $\card M$ is even, so $\card M\ge32$.
Take $T\subseteq M$ invariant with $\card T=32$, and put $T'=T\cup\{x\}\subseteq S$. Then $x\notin T$ (points of $T$ are
not fixed), so $\card{T'}=33$. $T'$ is invariant: $\sigma(x)=x$, and $T$ is invariant. The point $x\in T'$ is fixed, and
every fixed $y\in T'$ equals $x$, because no point of $T$ is fixed. $T'$ is in general position. \lean{MirrorOneAxis33}
applied to $T'$, $a$, $u$ gives \texttt{HasConvexNGon 7 T'}, hence \texttt{HasConvexNGon 7 S}.
\end{proof}

\subsection{The reduction theorems}\label{sr:ssec:main}

\begin{theorem}[Main reduction; \leanref{ErdosSzekeres.SymReduction.heptagon_of_nontrivial_isometry}{SYMLEAN/SymmetricReduction.lean:997}; conditional on five case classes]\label{sr:thm:main}
Assume the propositions \lean{HalfTurn34}, \lean{Threefold33}, \lean{Fivefold35}, \lean{MirrorOneAxis33} and
\lean{MirrorNoAxis34} of Definition~\ref{sr:def:classes} (hypotheses; none is proved in this section). Let $S$ be a
finite set of points and $g\colon\R^2\to\R^2$ a map. If $S$ is in general position, $33\le\card S$, $g$ is a Euclidean
isometry, $g\ne\mathrm{id}$ and $\operatorname{Invariant}(g,S)$, then \texttt{HasConvexNGon 7 S}.
\end{theorem}

\begin{proof}
Since $\card S\ge33>1$, Theorem~\ref{sr:thm:structure}(b) applies.

\emph{Rotation: $g=\operatorname{crot}_{c,a}$ with $a\ne1$ and $a^m=1$, $m>0$.} Then
$\operatorname{Invariant}(\operatorname{crot}_{c,a},S)$, and for every $k$,
$\operatorname{Invariant}(\operatorname{crot}_{c,a^k},S)$ by Lemmas~\ref{sr:lem:iterate}(b) and~\ref{sr:lem:crot}(d).
By Lemma~\ref{sr:lem:root-cases} there are two cases.
\begin{itemize}
\item $a^k=-1$ for some $k$. Then $S$ is invariant under $\operatorname{crot}_{c,-1}=
  \operatorname{halfTurn}_{\operatorname{ofC}(c)}$ (Lemma~\ref{sr:lem:rot-half}(b) and
  $\operatorname{toC}(\operatorname{ofC}c)=c$), and Lemma~\ref{sr:lem:case-halfturn} gives the claim.
\item $a^k$ is a primitive $p$-th root of unity for a prime $p\ne2$. Then Lemma~\ref{sr:lem:odd-prime}, applied with
  $w=a^k$, gives the claim.
\end{itemize}

\emph{Reflection: $g=\operatorname{reflLine}_{A,u}$ with $u_0^2+u_1^2=1$.} Lemma~\ref{sr:lem:mirror} gives the claim.
\end{proof}

\begin{remark}[Status of Theorem~\ref{sr:thm:main}]\label{sr:rem:main}
The theorem is Lean-checked as an implication: the five classes are explicit hypotheses (Lean binders
\texttt{h2 h3 h5 hm1 hm0}), not axioms, and the axioms log is \dirn{SYMLEAN/logs/SymmetricReductionAxioms.out}. The
unconditional statement, Theorem~\ref{sr:thm:main} without its five hypotheses, is claimed in the project as
computer-assisted (Theorem~A of \dirn{WRITEUP/SYMMETRIC.md}, \S0, in its form for Euclidean isometries), on the
evidence summarized in Remark~\ref{sr:rem:records}; it is not Lean-checked as a whole. In this document the theorem is
established only as the stated implication. Theorem~\ref{mb:thm:main3} restates the theorem with the two mirror
classes discharged, so that only \lean{HalfTurn34}, \lean{Threefold33} and \lean{Fivefold35} remain hypotheses; it is
stated in \dirn{SYMLEAN/mirror/ReductionThreeClasses.lean}, for the byte-identical copies of the definitions described
in Remark~\ref{sr:rem:records}. The
symmetry $g$ must be a Euclidean isometry of the whole plane; a distance-preserving permutation of $S$ alone is not
treated. Affine symmetries are treated in Theorem~\ref{ar:thm:main}.
\end{remark}

\begin{theorem}[Variant; \leanref{ErdosSzekeres.SymReduction.halfTurn_of_symmetric_counterexample}{SYMLEAN/SymmetricReduction.lean:1024}; conditional on five case classes]\label{sr:thm:variant}
Assume the propositions \lean{Threefold33}, \lean{Fourfold32}, \lean{Fivefold35}, \lean{MirrorOneAxis33} and
\lean{MirrorNoAxis34} of Definition~\ref{sr:def:classes} (hypotheses; none is proved in this section). Let $S$ be a
finite set of points and $g\colon\R^2\to\R^2$ a map. Assume that $S$ is in general position, $33\le\card S$, $g$ is a
Euclidean isometry, $g\ne\mathrm{id}$, $\operatorname{Invariant}(g,S)$, and $\neg$\,\texttt{HasConvexNGon 7 S}. Then
there is $c\in\R^2$ such that $g=\operatorname{halfTurn}_c$ (as maps of $\R^2$), $c\notin S$, $2\mid\card S$ and
$34\le\card S$.
\end{theorem}

\begin{proof}
By Theorem~\ref{sr:thm:structure}(b), as $\card S>1$, $g$ is a rotation or a reflection.

\emph{Rotation: $g=\operatorname{crot}_{c,a}$ with $a\ne1$ and $a^m=1$, $m>0$.} By Lemma~\ref{sr:lem:root-cases2} there
are three cases.
\begin{itemize}
\item $a=-1$. Then $g=\operatorname{crot}_{c,-1}=\operatorname{halfTurn}_{\operatorname{ofC}(c)}$
  (Lemma~\ref{sr:lem:rot-half}(b) and $\operatorname{toC}(\operatorname{ofC}c)=c$). By Lemma~\ref{sr:lem:halfturn-facts}, $\operatorname{ofC}(c)\notin S$ and
  $2\mid\card S$; with $\card S\ge33$ this gives $\card S\ge34$. So the point $\operatorname{ofC}(c)$ has all the
  required properties.
\item $a^k$ is a primitive $4$th root of unity. $S$ is invariant under $\operatorname{crot}_{c,a^k}$
  (Lemmas~\ref{sr:lem:iterate}(b) and~\ref{sr:lem:crot}(d)), and Lemma~\ref{sr:lem:fourfold} gives
  \texttt{HasConvexNGon 7 S}, contradicting the hypothesis.
\item $a^k$ is a primitive $p$-th root of unity for a prime $p\ne2$. Likewise Lemma~\ref{sr:lem:odd-prime} gives a
  contradiction.
\end{itemize}

\emph{Reflection.} Lemma~\ref{sr:lem:mirror} gives \texttt{HasConvexNGon 7 S}, a contradiction.
\end{proof}

\begin{remark}[Scope of Theorem~\ref{sr:thm:variant}]\label{sr:rem:variant}
The theorem concerns one symmetry at a time and Euclidean isometries only. Corollary~3 of
\dirn{WRITEUP/SYMMETRIC.md} has the same per-symmetry form (every symmetry of such a set is a half-turn about a point
outside the set), with the case theorems and the fourfold exclusion as assumptions, but there a symmetry may also be
an affine bijection; the theorem is the isometry case of that corollary. A group form, stated as the corollary labelled
\texttt{cor:only} in \dirn{WRITEUP/note/symmetric_es7.tex}, is Lean-checked for Euclidean isometries, with the same five
classes as hypotheses, in the review file \dirn{SYMLEAN/review/VariantChecks.lean} (\lean{VariantChecks.group_form},
line 1168): for $S$ as in the theorem, the Euclidean isometries that map $S$ into itself are the identity and at most
one half-turn, whose centre is not in $S$. The evidence is the log \dirn{SYMLEAN/logs/VariantChecks.out}, which prints
the axioms \lean{propext}, \lean{Classical.choice} and \lean{Quot.sound} for \lean{VariantChecks.group_form}, and the
row \texttt{review/VariantChecks.lean} of \dirn{SYMLEAN/logs/runs.tsv}, which records exit code~$0$ for the file with
SHA-256 prefix \texttt{4b823bb6}; the file at the commit used here has the prefix \texttt{4b823bb6cf85}. The group form
follows from the theorem. Let $g$ and $g'$ be Euclidean isometries other than the identity that map $S$ into itself.
The theorem gives $g=\operatorname{halfTurn}_c$ and $g'=\operatorname{halfTurn}_d$ with $c,d\notin S$. If $c\ne d$,
then $g\circ g'$ is the translation $p\mapsto p+2(c-d)$. It is a Euclidean isometry other than the identity that maps
$S$ into itself, so by the theorem it equals $\operatorname{halfTurn}_e$ for some $e\in\R^2$; evaluating at $p=e$ gives
$e+2(c-d)=e$, so $c=d$, a contradiction. The group form is not restated in this document as a numbered statement, and
no statement of this document uses it. The variant does not use \lean{HalfTurn34}; it uses \lean{Fourfold32}, which
rests on \dirn{C4CUT} (Remark~\ref{sr:rem:records}).
\end{remark}

\begin{definition}[The \dirn{H3} mirror statement; \leanref{ErdosSzekeres.SymReduction.MirrorExclusionLineCopy}{SYMLEAN/SymmetricReduction.lean:1052}]\label{sr:def:exclusion-line}
\lean{MirrorExclusionLineCopy} is the proposition: for every finite $S\subseteq\R^2$ and all $a,u\in\R^2$, if
$u_0^2+u_1^2=1$, $S$ is in general position, $\card S=33$, and $\operatorname{reflLine}_{a,u}(q)\in S$ for every $q\in
S$, then \texttt{HasConvexNGon 7 S}.
\end{definition}

\begin{lemma}[Bridge; \leanref{ErdosSzekeres.SymReduction.mirrorOneAxis33_of_exclusionLine}{SYMLEAN/SymmetricReduction.lean:1056}]\label{sr:lem:mirror-line}
\lean{MirrorExclusionLineCopy} implies \lean{MirrorOneAxis33}.
\end{lemma}

\begin{proof}
Let $S$, $a$, $u$ satisfy the hypotheses of \lean{MirrorOneAxis33}. The condition
$\operatorname{Invariant}(\operatorname{reflLine}_{a,u},S)$ is, by definition, the condition
$\operatorname{reflLine}_{a,u}(q)\in S$ for every $q\in S$. So all hypotheses of \lean{MirrorExclusionLineCopy} hold
(the hypothesis about the fixed point is not used), and it gives \texttt{HasConvexNGon 7 S}.
\end{proof}

\begin{remark}[About the bridge]\label{sr:rem:mirror-line}
\lean{MirrorExclusionLineCopy} has the same text as \lean{Mir.MirrorExclusionLine}
(\dirn{H3/MirrorEncoding/Line.lean:94}), and \lean{reflLine} has the same text as \lean{Mir.reflLine}. They are
different Lean constants of separately compiled developments, and no Lean theorem in \dirn{SYMLEAN} identifies them;
\dirn{SYMLEAN/README.md} \S5 reports a separate check of both identities by \texttt{rfl}, whose file is not kept in
\dirn{SYMLEAN}. In \dirn{H3}, \lean{Mir.mirrorExclusionLine_of_unsat} and
\lean{Mir.mirrorExclusionLine_of_cubes} (\dirn{H3/MirrorEncoding/Line.lean:139} and \dirn{:147}) derive
\lean{Mir.MirrorExclusionLine} from hypotheses saying that certain CNF formulas are unsatisfiable (the second also from
two hypotheses that the Boolean cover check \texttt{coverCheck} of a cube list returns \texttt{true}). In \dirn{V2-L},
\lean{V2L.mirrorExclusionLine_final} (\dirn{V2-L/V2L/FinalTree.lean:28}; \dirn{V2-L/README.md} \S9.1 and \S9.6)
derives it from hypotheses that all carry recorded verdicts. It is Lean-checked with four hypotheses, \texttt{h1},
\texttt{hroot}, \texttt{hup} and \texttt{hleaves}, which together assert $1{,}198+1+121+74{,}000$ facts of the form
``this CNF formula is unsatisfiable''; the $1+121$ facts of \texttt{hroot} and \texttt{hup} are the cover obligations
of the top two levels of the hull-pair case split of sub-case (ii) (Remark~\ref{sr:rem:records};
\dirn{WRITEUP/MIRROR.md} \S4). Each fact has a recorded \texttt{cake\_lpr} verdict on a file whose proof was
not retained, and the step from a verdict on a file to the Lean hypothesis is made outside Lean
(\dirn{WRITEUP/MIRROR.md} \S6). So
Lemma~\ref{sr:lem:mirror-line} supplies \lean{MirrorOneAxis33} only conditionally on those CNF facts and on the
identification of the two texts; the Lean proof without SAT input is in \dirn{SYMLEAN/mirror}
(Remark~\ref{sr:rem:records}). The converse implication also holds, and is Lean-checked in the review file
\dirn{SYMLEAN/review/ReviewSanity.lean} (\lean{Review.exclusionLine_of_mirrorOneAxis33}, line 1593, and the
equivalence \lean{Review.mirrorOneAxis33_iff_exclusionLine}, line 1629, whose axioms the log
\dirn{SYMLEAN/logs/ReviewSanity.out} prints): for $33$ points the fixed-point hypothesis of \lean{MirrorOneAxis33} is
automatic. Indeed, let $S$, $a$, $u$ satisfy the hypotheses of \lean{MirrorExclusionLineCopy}. In the notation of the
proof of Lemma~\ref{sr:lem:mirror}, whose first two steps use only these hypotheses, $\card F\le2$, $\card M$ is even
and $\card F+\card M=33$, so $\card F=1$. The point of $F$ is then the only point of $S$ fixed by
$\operatorname{reflLine}_{a,u}$, and \lean{MirrorOneAxis33} gives \texttt{HasConvexNGon 7 S}. Neither review theorem
is restated in this document as a numbered statement, and no statement of this document uses them.
\end{remark}

\begin{remark}[Hypotheses of the classes that follow from the others]\label{sr:rem:redundant}
Some hypotheses of the classes are implied by the others; this is not used anywhere.
\begin{itemize}
\item In \lean{HalfTurn34}, $c\notin S$ follows from general position and invariance (Lemma~\ref{sr:lem:center}).
\item In \lean{Fourfold32}, $c\notin S$ follows as well: by Lemma~\ref{sr:lem:rot-half}(a) the rotation is
  $f=\operatorname{crot}_{\hat c,\zeta}$ with $\zeta=\exp((2\pi/4)i)$, a primitive $4$th root of unity
  (Lemma~\ref{sr:lem:rot345}), so $\zeta^2=-1$ by (P2) and (P3). By Lemmas~\ref{sr:lem:crot}(d)
  and~\ref{sr:lem:rot-half}(b), $f^{[2]}=\operatorname{crot}_{\hat c,-1}=\operatorname{halfTurn}_c$; $S$ is invariant under
  it (Lemma~\ref{sr:lem:iterate}(b)), and Lemma~\ref{sr:lem:center} applies.
\item In \lean{Threefold33} and \lean{Fivefold35}, $c\notin S$ follows from the size. Let $p\in\{3,5\}$ and
  $w=\exp((2\pi/p)i)$, so that $w^p=1$ and $w\ne1$ (Lemma~\ref{sr:lem:rot345} and (P4)), and the rotation is
  $\operatorname{crot}_{\hat c,w}$ (Lemma~\ref{sr:lem:rot-half}(a)). If $c\in S$, then by Lemma~\ref{sr:lem:rot-erase}
  (with $\operatorname{ofC}(\hat c)=c$), $p$ divides $\card{(S\setminus\{c\})}$, which is $32$ for $p=3$ and $34$ for
  $p=5$; but $3\nmid32$ and $5\nmid34$.
\item In \lean{MirrorOneAxis33}, the fixed-point hypothesis follows from $\card S=33$ (Remark~\ref{sr:rem:mirror-line}).
  In \lean{MirrorNoAxis34} the hypothesis that no point is fixed is not implied in this way: for $34$ points the same
  argument gives only $\card F\in\{0,2\}$.
\end{itemize}
Dropping one of the implied hypotheses above gives an equivalent class. In general, adding a hypothesis to a class can
only make it weaker (the class without the hypothesis implies the class with it), that is, easier to supply.
\end{remark}

\paragraph{Coverage.}
All $72$ declarations of \dirn{SYMLEAN/SymmetricReduction.lean} are covered.
\begin{itemize}
\item Definitions: \lean{eDist2}, \lean{IsEucIsometry} (Definition~\ref{sr:def:iso}); \lean{Invariant}
  (Definition~\ref{sr:def:inv}); \lean{toC}, \lean{ofC} (Definition~\ref{sr:def:complex}); \lean{crot}, \lean{rotAbout},
  \lean{halfTurn}, \lean{reflLine} (Definition~\ref{sr:def:maps}); \lean{orbitFS} (Definition~\ref{sr:def:orbitFS}); the
  six classes \lean{HalfTurn34}, \lean{Threefold33}, \lean{Fivefold35}, \lean{MirrorOneAxis33}, \lean{MirrorNoAxis34},
  \lean{Fourfold32} (Definition~\ref{sr:def:classes}); \lean{MirrorExclusionLineCopy}
  (Definition~\ref{sr:def:exclusion-line}).
\item Finite dynamics: \lean{iterate_mem}, \lean{Invariant.iterate} (Lemma~\ref{sr:lem:iterate});
  \lean{mem_orbitFS_self}, \lean{orbitFS_subset}, \lean{orbitFS_invariant}, \lean{card_orbitFS},
  \lean{card_orbitFS_of_minimalPeriod} (Lemma~\ref{sr:lem:orbitFS}); \lean{orbit_select_minimalPeriod},
  \lean{orbit_select} (Lemma~\ref{sr:lem:orbit}).
\item Coordinates and maps: \lean{toC_re}, \lean{toC_im}, \lean{ofC_zero}, \lean{ofC_one}, \lean{toC_ofC},
  \lean{ofC_toC}, \lean{toC_injective}, \lean{eDist2_eq_normSq}, \lean{eDist2_pos_of_ne}, \lean{point_ext}
  (Lemma~\ref{sr:lem:complex}); \lean{toC_crot}, \lean{crot_one}, \lean{crot_comp}, \lean{crot_iterate},
  \lean{crot_eq_self_iff}, \lean{crot_injective}, \lean{eDist2_crot_center} (Lemma~\ref{sr:lem:crot});
  \lean{rotAbout_eq_crot}, \lean{halfTurn_eq_crot} (Lemma~\ref{sr:lem:rot-half}); \lean{toC_reflLine},
  \lean{reflLine_involutive}, \lean{cross_of_reflLine_fixed}, \lean{orient_eq_zero_of_cross}
  (Lemma~\ref{sr:lem:refl}).
\item Isometries: \lean{IsEucIsometry.injective}, \lean{isometry_toC_conj} (Lemma~\ref{sr:lem:toC-iso});
  \lean{isometry_complex_cases}, \lean{isometry_structure} (Theorem~\ref{sr:thm:structure}).
\item Geometry: \lean{convexPosition_of_concyclic} (Lemma~\ref{sr:lem:concyclic}), \lean{hasConvexNGon_of_concyclic}
  (Corollary~\ref{sr:cor:concyclic}), \lean{center_not_mem_of_halfTurn} (Lemma~\ref{sr:lem:center}); also the library
  lemma \lean{Suk.convexPosition_of_strict_supports} (Lemma~\ref{sr:lem:supports}). This is the lemma titled ``Strict
  supporting lines'' in the foundations section; the short proof is repeated here.
\item Roots of unity: \lean{root_cases} (Lemma~\ref{sr:lem:root-cases}); \lean{prime_cases}
  (Lemma~\ref{sr:lem:prime-cases}); \lean{root_cases2} (Lemma~\ref{sr:lem:root-cases2}); and
  \lean{isPrimitiveRoot_rot3}, \lean{isPrimitiveRoot_rot5}, \lean{isPrimitiveRoot_rot4} (Lemma~\ref{sr:lem:rot345}).
\item Cases: \lean{halfTurn_facts} (Lemma~\ref{sr:lem:halfturn-facts}); \lean{case_halfTurn}
  (Lemma~\ref{sr:lem:case-halfturn}, conditional on \lean{HalfTurn34}); \lean{rot_erase_facts}
  (Lemma~\ref{sr:lem:rot-erase}); \lean{case_rot_big} (Lemma~\ref{sr:lem:rot-big}); \lean{case_odd_prime}
  (Lemma~\ref{sr:lem:odd-prime}, conditional on \lean{Threefold33} and \lean{Fivefold35}); \lean{case_fourfold}
  (Lemma~\ref{sr:lem:fourfold}, conditional on \lean{Fourfold32}); \lean{case_mirror} (Lemma~\ref{sr:lem:mirror},
  conditional on \lean{MirrorOneAxis33} and \lean{MirrorNoAxis34}).
\item Main results: \lean{heptagon_of_nontrivial_isometry} (Theorem~\ref{sr:thm:main}, conditional on
  \lean{HalfTurn34}, \lean{Threefold33}, \lean{Fivefold35}, \lean{MirrorOneAxis33} and \lean{MirrorNoAxis34});
  \lean{halfTurn_of_symmetric_counterexample} (Theorem~\ref{sr:thm:variant}, conditional on the same classes with
  \lean{Fourfold32} in place of \lean{HalfTurn34});
  \lean{mirrorOneAxis33_of_exclusionLine} (Lemma~\ref{sr:lem:mirror-line}, an implication between two propositions,
  so conditional on its hypothesis \lean{MirrorExclusionLineCopy}).
\end{itemize}
The local statement \lean{pigeon} in the proof of \lean{isometry_structure} is Lemma~\ref{sr:lem:pigeon}. The Mathlib
facts on minimal periods and on roots of unity that the Lean proofs use are restated and proved in
Lemmas~\ref{sr:lem:minper} and~\ref{sr:lem:prim}; the Mathlib results on isometries that the Lean proof uses (the
construction \lean{Isometry.affineIsometryOfStrictConvexSpace} and the theorem \lean{linear_isometry_complex}) are
replaced by the self-contained proof of Theorem~\ref{sr:thm:structure}(a); other standard facts (the pigeonhole
principle, least prime factors, square roots of complex numbers, elementary facts on finite sets) are stated where they
are used. Not covered,
deliberately:
\dirn{SYMLEAN/SymmetricReductionAxioms.lean} (a verbatim copy with two \texttt{\#print axioms} lines); lines 1--1058 of
\dirn{SYMLEAN/AffineReduction.lean} (a verbatim copy; its affine part belongs to the section on affine symmetries);
the review files \dirn{SYMLEAN/review/ReviewSanity.lean}, \dirn{SYMLEAN/review/VariantChecks.lean} and
\dirn{SYMLEAN/review/SupMetric.lean} (review and audit files, cited only in the paragraph on axioms and rigor and in
Remarks~\ref{sr:rem:sup}, \ref{sr:rem:variant} and~\ref{sr:rem:mirror-line}); and \dirn{SYMLEAN/mirror} (the section
on mirror bounds). Two results of the review files go beyond audit checks and are not restated as numbered
statements: the group form \lean{VariantChecks.group_form} (\dirn{SYMLEAN/review/VariantChecks.lean:1168}, conditional
on the five classes of Theorem~\ref{sr:thm:variant}), and the converse bridge
\lean{Review.exclusionLine_of_mirrorOneAxis33} (\dirn{SYMLEAN/review/ReviewSanity.lean:1593}, an implication between
two propositions, so conditional on its hypothesis \lean{MirrorOneAxis33}) with the equivalence
\lean{Review.mirrorOneAxis33_iff_exclusionLine} of the two propositions (line~1629 of the same file), which asserts
neither of them. The reason is that they lie in review files and that no statement of this document uses them.
Remarks~\ref{sr:rem:variant} and~\ref{sr:rem:mirror-line} state them, name
their logs and give short written arguments.
